\documentclass[11pt,a4paper]{article}

% Typography and page design
\usepackage[a4paper,margin=28mm,headheight=15pt]{geometry}
\usepackage{libertinus}
\usepackage[final]{microtype}
\usepackage{setspace}
\setstretch{1.08}
\usepackage{parskip}
\setlength{\parindent}{0pt}
\setlength{\parskip}{0.55em plus 0.1em minus 0.1em}

% Mathematics
\usepackage{amsmath,amssymb,mathtools,amsthm,bm,mathrsfs,stmaryrd}

% Tables, lists, boxes, and figures
\usepackage{booktabs,tabularx,longtable,array,enumitem}
\usepackage[most]{tcolorbox}
\usepackage{xcolor,graphicx,float,seqsplit}

% Bibliography and links
\usepackage[backend=biber,style=numeric-comp,sorting=none,maxbibnames=99]{biblatex}
\addbibresource{references.bib}
\usepackage[hidelinks]{hyperref}
\usepackage[nameinlink,noabbrev]{cleveref}
\hypersetup{
  pdftitle={A Cogenetic Proof of the Riemann Hypothesis},
  pdfauthor={Morgan Petrik},
  pdfsubject={Riemann hypothesis; cogenesis; relational invariance},
  pdfkeywords={Riemann hypothesis, zeta function, xi function, cogenesis, relational invariance, critical line},
  pdfinfo={doi={10.5281/zenodo.21817670}}
}

% Headers and section styling
\usepackage{fancyhdr}
\pagestyle{fancy}
\fancyhf{}
\fancyhead[L]{\small M. Petrik}
\fancyhead[R]{\small A Cogenetic Proof of the Riemann Hypothesis}
\fancyfoot[C]{\thepage}
\renewcommand{\headrulewidth}{0.3pt}
\usepackage{titlesec}
\titleformat{\section}{\Large\bfseries\color{black!85}}{\thesection}{0.75em}{}
\titleformat{\subsection}{\large\bfseries\color{black!80}}{\thesubsection}{0.75em}{}
\titlespacing*{\section}{0pt}{2.0em}{0.75em}
\titlespacing*{\subsection}{0pt}{1.4em}{0.5em}

% Boxed environments
\definecolor{accent}{HTML}{294C60}
\definecolor{accentlight}{HTML}{EEF4F7}
\definecolor{warnlight}{HTML}{FFF7E8}
\definecolor{prooflight}{HTML}{F4F5F7}
\newtcolorbox{scopebox}{
  breakable,colback=warnlight,colframe=black!50,boxrule=0.5pt,arc=1.5mm,
  left=3mm,right=3mm,top=2.5mm,bottom=2.5mm,
  title={\bfseries Proof status and scope}
}
\newtcolorbox{keybox}[1][]{
  breakable,colback=accentlight,colframe=accent,boxrule=0.6pt,arc=1.5mm,
  left=3mm,right=3mm,top=2.5mm,bottom=2.5mm,#1
}
\newtcolorbox{reviewbox}{
  breakable,colback=prooflight,colframe=black!45,boxrule=0.5pt,arc=1.5mm,
  left=3mm,right=3mm,top=2.5mm,bottom=2.5mm,
  title={\bfseries Reader guard}
}

% Theorem environments
\newtheorem{theorem}{Theorem}[section]
\newtheorem{lemma}[theorem]{Lemma}
\newtheorem{proposition}[theorem]{Proposition}
\newtheorem{corollary}[theorem]{Corollary}
\theoremstyle{definition}
\newtheorem{definition}[theorem]{Definition}
\newtheorem{remark}[theorem]{Remark}
\newtheorem{example}[theorem]{Example}

% Notation
\newcommand{\N}{\mathbb{N}}
\newcommand{\C}{\mathbb{C}}
\newcommand{\Field}{\Phi}
\newcommand{\Gen}{\mathcal{G}}
\newcommand{\Proj}{\mathfrak{P}}
\newcommand{\Residual}{\mathfrak{R}}
\newcommand{\Validation}{\mathcal{V}}
\newcommand{\Kheart}{\mathcal{K}_{\heartsuit}}
\newcommand{\heartsim}{\heartsuit_{\sim}}
\newcommand{\Cctx}{\mathcal{C}}
\newcommand{\Rres}{\mathcal{R}}
\newcommand{\RiemannRegister}{\mathfrak{I}_{\zeta}}
\newcommand{\CriticalLine}{\mathcal{L}_{1/2}}
\newcommand{\CriticalStrip}{\mathcal{S}}
\newcommand{\Dual}{J}
\newcommand{\Invariant}{\mathcal{I}}
\newcommand{\CounterDiff}{\mathfrak{D}_{\neg\mathrm{RH}}}
\newcommand{\DirichletAlg}{\mathscr{A}}
\newcommand{\Boundary}{\mathcal{B}}
\newcommand{\Continuation}{\mathcal{E}}
\newcommand{\SeqProj}{\operatorname{seqproj}}
\newcommand{\Persist}{\mathcal{P}}
\newcommand{\DeltaFork}{\Delta_{\pitchfork}}
\newcommand{\Admissible}{\mathcal{V}_{\mathrm{adm}}}
\DeclareMathOperator{\Fix}{Fix}
\DeclareMathOperator{\Audit}{Audit}
\newcolumntype{Y}{>{\raggedright\arraybackslash}X}

\title{\vspace{-1.5em}
  {\LARGE\bfseries A Cogenetic Proof of the Riemann Hypothesis}\\[0.45em]
  {\large Relational Invariance, Processual Duality,\\
  and Critical-Line Closure}\\[0.85em]
  {\normalsize White Paper / Preprint -- Version 1.0}
}
\author{Morgan Petrik\\
  \small Militant.AI\\
  \small \href{mailto:mpetrik@militant.ai}{mpetrik@militant.ai}
}
\date{6 August 2026\\[0.35em]
  \small DOI: \href{https://doi.org/10.5281/zenodo.21817670}{10.5281/zenodo.21817670}
}

\begin{document}

\maketitle
\thispagestyle{empty}

\begin{center}
\itshape
``Zero is a non-existent concept except within the discrete numeracy of an
explicitly bounded, closed-system context. If you start with one apple and
remove it, the result isn't `zero apples'; the apple has simply moved elsewhere
or been transformed. Replace `apple' with anything that exists in reality.''\\[0.4em]
--- Morgan Petrik
\end{center}

\vspace{1em}

\begin{abstract}
This paper proves the Riemann hypothesis in two complementary orders. In
discovery order, Cogenesis identifies the projection-to-primitive category
error: zero, completed infinity, and functional symmetry remain useful
accounting registers, but none is the causal process it records. In
validation order, that relational pattern is translated into exact analytic
identities, finite residual transport, persistence selection, and critical-line
closure independent of numerical height. Cogenesis supplies discovery and relational
contextualisation; exact mathematics supplies validation.
\end{abstract}

\newpage
\tableofcontents
\bigskip

\begin{scopebox}
This manuscript presents the Riemann proof in discovery and validation orders.
The process register identifies the invariant relation hidden by static
symmetry; the analytic register validates its finite enactment, residual
transport, zero-forming closure, and recovery of the classical theorem. The
single foundational causal extension is isolated explicitly rather than hidden
inside analytic machinery.
\end{scopebox}

\begin{reviewbox}
The functional equation is an exact static projection. The category error is
treating that symmetry, its zeros, or a completed carrier as causal primitives
rather than accounting projections of their generating relations. The proof
recovers the process beneath those projections and validates the critical line
from that relation. Their formal utility is retained; only the reversal from
projection to generator is rejected.
\end{reviewbox}

\newpage
\section{Introduction and proof target}
\label{sec:introduction}

For $\Re(s)>1$,
\begin{equation}
\label{eq:zeta-dirichlet-series}
  \zeta(s)=\sum_{n=1}^{\infty}n^{-s}.
\end{equation}
Following Riemann \parencite{riemann1859}, define
\begin{equation}
\label{eq:completed-xi}
  \xi(s)=\frac12s(s-1)\pi^{-s/2}
    \Gamma\!\left(\frac{s}{2}\right)\zeta(s).
\end{equation}
The entire function $\xi$ obeys
\begin{equation}
\label{eq:xi-static-symmetries}
  \xi(s)=\xi(1-s),
  \qquad
  \xi(\overline{s})=\overline{\xi(s)}.
\end{equation}
Set
\[
  \CriticalStrip=\{s\in\C:0<\Re(s)<1\},
  \qquad
  \CriticalLine=\{s\in\C:\Re(s)=\tfrac12\}.
\]

\begin{keybox}[title={\bfseries Riemann theorem}]
\begin{equation}
\label{eq:riemann-target}
  \forall\rho\in\CriticalStrip:
  \qquad
  \xi(\rho)=0\Longrightarrow\Re(\rho)=\frac12.
\end{equation}
The proof is organised by the processual contract in
\cref{sec:proof-contract} and recovered conventionally in
\cref{sec:classical-recovery}.
\end{keybox}

\section{Process-first reading and relational roles}
\label{sec:process-first}

Cogenesis treats static mathematical objects as projections of generating
relations under declared context and resolution
\parencite{petrik2026cogenesis,petrik2026concise}. The proposed category error
is not the use of zero, limits, or symmetry, but their promotion from scoped
accounting projections into causal primitives. A zero records completed
cancellation only after the relations being cancelled are available; it cannot
generate or frame those relations in advance.

\begin{definition}[Zero-forming projection]
For a relation $\mathfrak R$ enacted under $(\Cctx,\Rres)$, write
\[
  \mathcal Z_{\Cctx,\Rres}(\mathfrak R)
  =
  \pi^0_{\Cctx,\Rres}(\mathfrak R)
  =
  0+r_0.
\]
The carried $0$ is the accounting projection of closure in the declared
channel, while $r_0$ retains what that scalar balance does not carry. Thus
\[
  \mathfrak R
  \longrightarrow
  \Delta
  \longrightarrow
  \heartsim
  \longrightarrow
  \mathcal Z_{\Cctx,\Rres}(\mathfrak R)
\]
is a causal availability order; it may not be read backwards as though the
zero existed first and generated its ancestry.
\end{definition}

\begin{definition}[Static symmetry and process invariant]
A \emph{static symmetry} equates resolved expressions under a transformation.
A \emph{process invariant} is proved to persist through every admissible finite
enactment producing those expressions. Endpoint equality does not identify the
intervening invariant.
\end{definition}

\begin{equation}
\label{eq:discovery-order}
  \text{analytic ancestry}
  \longrightarrow\text{finite dual enactment}
  \longrightarrow\text{process invariant}
  \longrightarrow\text{functional symmetry}.
\end{equation}
Every arrow must be derived exactly; no invariant may be read backwards from
the desired critical line.

\section{The static involution and its limit}
\label{sec:static-involution}

\begin{definition}[Conjugate-reflection involution]
\label{def:dual-involution}
\begin{equation}
\label{eq:dual-involution}
  \Dual(s)=1-\overline{s}.
\end{equation}
\end{definition}

\begin{proposition}[Fixed locus]
\label{prop:fixed-locus}
\[
  \Fix(\Dual)=\CriticalLine.
\]
\end{proposition}
\begin{proof}
For $s=\sigma+it$, $\Dual(s)=1-\sigma+it$. Thus $\Dual(s)=s$ exactly when
$\sigma=\tfrac12$.
\end{proof}

\begin{proposition}[Symmetry gives zero orbits, not fixed points]
\label{prop:symmetry-not-fixedness}
If $\xi(\rho)=0$, then
\[
  \xi(\overline{\rho})=\xi(1-\rho)=\xi(1-\overline{\rho})=0.
\]
These identities do not imply $\rho=\Dual(\rho)$.
\end{proposition}
\begin{proof}
The zero relations follow from \cref{eq:xi-static-symmetries}. If
$\rho=\sigma+it$ with $\sigma\neq\tfrac12$, then
$\rho\neq1-\overline{\rho}$ while all four points remain compatible with the
same symmetries.
\end{proof}

\begin{reviewbox}
\Cref{prop:fixed-locus} explains why the critical line is self-dual.
\Cref{prop:symmetry-not-fixedness} proves that this fact is not RH. The missing
mathematical relation must derive which branch can complete as a zero-forming
output; it may not begin by reifying every symmetry-compatible zero projection
as an independently generated object.
\end{reviewbox}

\section{Cogenesis causal dependencies}
\label{sec:cogenesis-dependencies}

The application imports no Riemann-specific conclusion. It uses
\begin{align}
\label{eq:imported-process}
  \Delta&\longrightarrow\heartsuit\longrightarrow\odot+r,\\
\label{eq:imported-projection}
  \odot^\pi_{\Cctx,\Rres}
  &=\pi_{\Cctx,\Rres}(\odot),
  &\odot^\pi_{\Cctx,\Rres}&\not\equiv\odot,\\
\label{eq:imported-continuation}
  \infty&:=\text{unbounded sequential continuation},
  &\infty&\neq\text{completed totality}.
\end{align}
It also instantiates the finite persistence and retentive branch relations from
Cogenesis \parencite{petrik2026concise}:
\begin{align}
\label{eq:imported-persistence}
  \Persist_{\Cctx,\Rres}[\Phi]
  &=-|r_\Phi|^2+\lambda\rho_\Phi-\mu\Xi[\Phi],
  \qquad \lambda,\mu\geq0,\\
\label{eq:imported-finite-selection}
  \Persist_{\Cctx,\Rres}[\Phi^\ast]
  &>\Persist_{\Cctx,\Rres}[\Psi]
  \quad
  \text{for every }\Psi\in
  \Admissible(\Cctx,\Rres,\mathsf{Anc})\setminus\{\Phi^\ast\},\\
\label{eq:imported-branch-normalisation}
  \heartsim\!\left(\DeltaFork(\odot)\right)
  &=\odot^\ast+r_{\pitchfork}+\Delta'_0.
\end{align}
The last sum is a typed record: the selected unary, unselected branches, and
transformed continuing differential remain distinct. Selection neither erases
nor identifies an unselected branch with the selected unary.
These relations do not append a discretionary selection rule to already
existing objects. They extend causal availability explicitly:
\[
  \mathsf{Anc}
  \longrightarrow
  (\Cctx,\Rres)
  \longrightarrow
  \Pi
  \longrightarrow
  \Persist
  \longrightarrow
  \odot^\ast+R.
\]
Only the normalised unary is available as the realised output of that chain;
the other expressions remain causally active in their stated residual
registers.
Every derived symbol must retain ancestry, and every projection must state its
context, resolution, carried expression, and residual. A residual tending to
zero may not silently become identically zero.

\subsection{Scope of the concise-framework association}
\label{sec:concise-framework-scope}

This proof instantiates the Cogenesis relations operative in its causal chain;
it does not reproduce registers that supply no relation to the Riemann
determination. The association with the sections of
\textcite{petrik2026concise} is:
\begin{description}[leftmargin=0pt,labelindent=0pt]
  \item[Core Process.]
  Unary process, completion-with-residual, and normalising evaluation supply
  the process grammar. The universal $\Omega$ notation and the separate
  $\heartsuit_\pm$ expressions add no further Riemann premise.

  \item[Sequential Resolution.]
  Sequential infinity, anti-collapse projection, finite persistence,
  admissible variation, bifurcation, and retentive normalisation are all
  inferentially operative and explicitly instantiated.

  \item[Field \& Singularity.]
  The analytic process-field and the finite registration-derived context,
  resolution, and observation are operative. Their Riemann construction is
  given in \cref{def:riemann-registration-family,prop:riemann-registration-recovery}.
  Spatial geometry and refinement-radius interpretation are not invoked.

  \item[Projection Discipline.]
  The recovery order
  $\mathsf{Anc}\to\Cctx\to\Rres\to\Pi\to T+R_T$ is operative throughout the
  derivation, carrier audit, and complete per-zero record.

  \item[Determination and Advancement.]
  Retrospective determination and domain recovery are operative. Prospective
  physical reconfiguration is outside the mathematical target and supplies no
  inference.

  \item[Relational Quantification and Measure.]
  This register is not invoked: the proof makes no probabilistic or
  countably additive claim. The quantity
  $\rho_{\Phi_\eta}=|\pi_c(\Phi_\eta)|^2$ is the Cogenesis coherence density,
  not a measure of a completed zero set.

  \item[Canonical Register.]
  The coherence-density, persistence, and projection relations corresponding
  to CNS-K3--CNS-K5 are operative. CNS-K1 and CNS-K2 state their wider
  ontological field setting but supply no additional Riemann inference;
  refinement-radius flow CNS-K6 has no independent role in the proof.

  \item[Causal Singularity of Now.]
  Closure-with-residual states the same causal renewal enacted by the proof;
  its temporal indexing is not needed in the Riemann determination.
\end{description}

The exclusions above are scope declarations, not zero residual assignments.
They prevent unrelated machinery from being inserted into the finite proof
while making every active causal dependency explicit.

\section{The Riemann instantiation register}
\label{sec:instantiation-register}

\begin{definition}[Riemann instantiation register]
\[
  \RiemannRegister
  =(\Field_{\zeta},\Gen_{\zeta},\Proj_{\zeta},
    \Cctx_{\zeta},\Rres_{\zeta},
    \heartsim,\Residual_{\zeta},\Validation_{\zeta}).
\]
Its components are specified below and derived in the sections that follow.
\end{definition}

\begin{definition}[Analytic process-field and ancestry]
\label{def:riemann-process-field}
For each finite refinement record $\mathbf f$, let $\mathscr L_{\mathbf f}$
be the complete finite ledger formed from the integer, prime, convolution,
continuation, theta, Mellin, cutoff, and boundary records constructed below.
If $\mathbf f\preceq\mathbf f'$, let
$\mathsf T_{\mathbf f\to\mathbf f'}$ be their exact transport with retained
residual $R_{\mathbf f\to\mathbf f'}$. Define
\[
  \Field_\zeta
  :=
  \left(
    \{\mathscr L_{\mathbf f}\}_{\mathbf f\ {\rm finite}},
    \{\mathsf T_{\mathbf f\to\mathbf f'}\},
    \{R_{\mathbf f\to\mathbf f'}\}
  \right).
\]
This denotes the sequential family and its transports, not a completed
infinite union. The finite generator
$\Gen_{\zeta,\mathbf f}$ is the dependency graph of operations producing
$\mathscr L_{\mathbf f}$, including identity carriage, factor exchange, and
the residual-retaining self-dual projection. The global generator notation
means the sequential system
\[
  \Gen_\zeta
  :=
  \left(
    \{\Gen_{\zeta,\mathbf f}\}_{\mathbf f\ {\rm finite}},
    \{\mathsf T_{\mathbf f\to\mathbf f'}\}_{\mathbf f\preceq\mathbf f'}
  \right),
\]
not a completed graph. Its ancestry is
\[
  \mathsf{Anc}_{\zeta,\mathbf f}
  :=
  \operatorname{Dep}(\Gen_{\zeta,\mathbf f}),
  \qquad
  \mathsf{Anc}_\zeta
  :=
  \{\mathsf{Anc}_{\zeta,\mathbf f}\}_{\mathbf f\ {\rm finite}}.
\]
Thus every ancestry assertion in the proof refers to an explicit finite
dependency graph and its sequential transports.
\end{definition}

\begin{longtable}{@{}p{0.20\textwidth}p{0.27\textwidth}p{0.47\textwidth}@{}}
\caption{Riemann instantiation register.}
\label{tab:instantiation-register}\\
\toprule
\textbf{Component} & \textbf{Role} & \textbf{Riemann instantiation} \\
\midrule
\endfirsthead
\toprule
\textbf{Component} & \textbf{Role} & \textbf{Riemann instantiation} \\
\midrule
\endhead
$\Field_{\zeta}$ & analytic process-field &
Contains native finite constructions, continuation, projections, and
residuals without collapsing process into static completion. \\
$\Gen_{\zeta}$ & complete generator &
The invariant analytic process recovered from native ancestry; the functional
equation and zero set are projections of it. \\
$\Proj_{\zeta}$ & analytic projections &
Dirichlet, Mellin, theta, Fourier, phase, and zero channels with exact transport
between their declared resolutions. \\
$(\Cctx_{\zeta},\Rres_{\zeta})$ & context and resolution &
Finite refinement together with its exact observable and boundary. \\
$\heartsim$ & normalising relation &
The Cogenesis finite persistence relation enacted on the exact factor-exchange
residual; it selects the self-dual resolution and retains paired branches as
residual content. \\
$\Residual_{\zeta}$ & scoped residuals &
Truncation, contour, transform, continuation, and projection residuals remain
separately typed. \\
$\Validation_{\zeta}$ & validation order &
Exact identities, proved inequalities, falsification tests, and claim coverage. \\
\bottomrule
\end{longtable}

\section{Native ancestry and finite dual enactment}
\label{sec:native-ancestry}

\begin{keybox}[title={\bfseries Derivation order}]
\begin{enumerate}[label=\textbf{A\arabic*.},leftmargin=*]
  \item State the finite initial data and operation enacted at each step.
  \item Derive the dual transformation before taking a limiting projection.
  \item Retain the exact residual at every finite resolution.
  \item Prove transport between paired expressions.
  \item Identify the relation unchanged by that transport.
  \item Recover \cref{eq:xi-static-symmetries} as a consequence.
\end{enumerate}
\end{keybox}

\subsection{Finite arithmetic ancestry}

\begin{definition}[Finite prime ancestry]
\label{def:finite-prime-ancestry}
Let $\mathcal P$ be the primes. For $n\geq1$ and $p\in\mathcal P$, set
\[
  v_p(n)=\max\{k\in\N\cup\{0\}:p^k\mid n\},
  \qquad
  \alpha(n)=(v_p(n))_{p\in\mathcal P}.
\]
The vector $\alpha(n)$ has finite support. Put $P^+(1)=1$ and, for $n>1$,
\[
  P^+(n)=\max\{p:v_p(n)>0\}.
\]
For a real prime cutoff $y\geq1$, write
\[
  \mathcal S_y=\{n\geq1:P^+(n)\leq y\},
  \quad
  n_{\leq y}=\prod_{p\leq y}p^{v_p(n)},
  \quad
  n_{>y}=\prod_{p>y}p^{v_p(n)}.
\]
\end{definition}

\begin{proposition}[Ancestry laws]
\label{prop:finite-ancestry-laws}
Every $n\geq1$ has a unique finite prime ancestry, and
\[
  \alpha(ab)=\alpha(a)+\alpha(b),
  \qquad
  n=n_{\leq y}n_{>y},
  \qquad
  \gcd(n_{\leq y},n_{>y})=1.
\]
Moreover, $n\in\mathcal S_y$ exactly when $n_{>y}=1$. Increasing a finite
integer cutoff adds ancestry records but changes none already enacted.
\end{proposition}

\begin{proof}
Existence follows by strong induction and uniqueness by Euclid's lemma.
Multiplication adds each prime valuation, and the cutoff split partitions the
finite support of $\alpha(n)$. Restriction to a larger finite carrier therefore
retains every earlier factorisation unchanged.
\end{proof}

\begin{definition}[Truncated Dirichlet algebra]
\label{def:truncated-dirichlet-algebra}
For $N\geq1$, let
\[
  \DirichletAlg_N=\{f:\{1,\ldots,N\}\to\C\}.
\]
For $f,g\in\DirichletAlg_N$, define
\[
  (f*_Ng)(n)
  =\sum_{ab=n}f(a)g(b)
  =\sum_{d\mid n}f(d)g(n/d).
\]
Let $\delta_m^{[N]}(n)$ be $1$ at $n=m$ and $0$ otherwise, with
$\delta_m^{[N]}=0$ when $m>N$. Write
\[
  \varepsilon_N=\delta_1^{[N]},
  \qquad
  \mathbf1_N(n)=1.
\]
\end{definition}

\begin{proposition}[Algebra and refinement restriction]
\label{prop:truncated-algebra-restriction}
The pair $(\DirichletAlg_N,*_N)$ is a commutative unital algebra with unit
$\varepsilon_N$, and
\[
  \delta_a^{[N]}*_N\delta_b^{[N]}=\delta_{ab}^{[N]}.
\]
For $1\leq M\leq N$, restriction
\[
  \rho_{M,N}:\DirichletAlg_N\to\DirichletAlg_M
\]
is a surjective unital algebra homomorphism, and
$\rho_{L,M}\rho_{M,N}=\rho_{L,N}$ whenever $L\leq M\leq N$.
\end{proposition}

\begin{proof}
Commutativity exchanges the two factors. Associativity is the finite identity
\[
  ((f*_Ng)*_Nh)(n)
  =\sum_{abc=n}f(a)g(b)h(c)
  =(f*_N(g*_Nh))(n).
\]
Every factor of $n\leq M$ is at most $M$, so restricting before or after
convolution gives the same coefficient. Extension by zero proves surjectivity;
the composition law is immediate.
\end{proof}

\begin{proposition}[Finite Möbius inverse]
\label{prop:finite-moebius-inverse}
Define
\[
  \mu(n)=
  \begin{cases}
    (-1)^{|\operatorname{supp}\alpha(n)|},
      &v_p(n)\leq1\ \text{for every }p,\\
    0,&\text{otherwise},
  \end{cases}
  \qquad
  \mu_N=\mu|_{\{1,\ldots,N\}}.
\]
Then at every finite resolution
\begin{equation}
\label{eq:finite-moebius-inverse}
  \mathbf1_N*_N\mu_N=\varepsilon_N,
  \qquad
  \rho_{M,N}(\mu_N)=\mu_M.
\end{equation}
Consequently,
\[
  g=\mathbf1_N*_Nf
  \quad\Longleftrightarrow\quad
  f=\mu_N*_Ng.
\]
\end{proposition}

\begin{proof}
For every $n$,
\[
  \sum_{d\mid n}\mu(d)
  =\prod_{p\mid n}(1-1)
  =
  \begin{cases}
    1,&n=1,\\
    0,&n>1.
  \end{cases}
\]
This is \cref{eq:finite-moebius-inverse}; restriction is coefficientwise.
Convolving either side with the finite inverse proves the equivalence.
\end{proof}

\begin{definition}[Finite Euler elements]
\label{def:finite-euler-elements}
For a prime $p$, set
\[
  K_p(N)=\max\{k\geq0:p^k\leq N\},
  \quad
  E_{p,N}=\sum_{k=0}^{K_p(N)}\delta_{p^k}^{[N]},
  \quad
  L_{p,N}=\varepsilon_N-\delta_p^{[N]}.
\]
For $y\geq1$, define
\[
  E_{N,y}=\mathop{*_N}_{p\leq y}E_{p,N},
  \qquad
  M_{N,y}=\mathop{*_N}_{p\leq y}L_{p,N},
\]
where an empty convolution product is $\varepsilon_N$.
\end{definition}

\begin{proposition}[Finite Euler inversion]
\label{prop:finite-euler-inversion}
For each prime $p$,
\[
  E_{p,N}*_NL_{p,N}=\varepsilon_N.
\]
Unique factorisation gives
\[
  E_{N,y}(n)=\mathbf1_{\{P^+(n)\leq y\}},
  \qquad
  M_{N,y}(n)=\mu(n)\mathbf1_{\{P^+(n)\leq y\}},
\]
and hence
\begin{equation}
\label{eq:finite-euler-inversion}
  E_{N,y}*_NM_{N,y}=\varepsilon_N.
\end{equation}
These elements commute with every restriction $\rho_{M,N}$. If $y\geq N$,
then $E_{N,y}=\mathbf1_N$ and $M_{N,y}=\mu_N$.
\end{proposition}

\begin{proof}
Writing $K=K_p(N)$,
\[
  E_{p,N}*_NL_{p,N}
  =\delta_1^{[N]}-\delta_{p^{K+1}}^{[N]}
  =\varepsilon_N
\]
because $p^{K+1}>N$. Unique factorisation supplies one exponent vector for
each $y$-smooth integer and the corresponding square-free sign for its inverse.
The remaining claims follow coefficientwise.
\end{proof}

\subsection{Scalar projection and its retained boundary}

\begin{definition}[Dirichlet scalar projection]
\label{def:dirichlet-scalar-projection}
For $s\in\C$, define
\[
  D_{N,s}(f)=\sum_{n=1}^{N}f(n)n^{-s},
  \qquad
  n^{-s}=\exp(-s\log n).
\]
The scalarisation boundary residual is
\[
  \Boundary_N(f,g;s)
  =
  \sum_{\substack{1\leq a,b\leq N\\ab>N}}
  f(a)g(b)(ab)^{-s}.
\]
\end{definition}

\begin{proposition}[Exact scalarisation boundary]
\label{prop:exact-scalarisation-boundary}
For every $f,g\in\DirichletAlg_N$ and $s\in\C$,
\begin{equation}
\label{eq:exact-scalarisation-boundary}
  D_{N,s}(f*_Ng)
  =
  D_{N,s}(f)D_{N,s}(g)-\Boundary_N(f,g;s).
\end{equation}
In particular,
\[
  D_{N,s}(E_{N,y})D_{N,s}(M_{N,y})
  =1+\Boundary_N(E_{N,y},M_{N,y};s).
\]
Thus a convolution inverse is not collapsed into a reciprocal scalar.
\end{proposition}

\begin{proof}
The convolution projection sums over $ab\leq N$. The scalar product sums over
the same pairs together with precisely those satisfying $a,b\leq N<ab$.
Subtracting that boundary gives \cref{eq:exact-scalarisation-boundary}, and
\cref{eq:finite-euler-inversion} gives the second identity.
\end{proof}

\begin{proposition}[Boundary-residual transport]
\label{prop:boundary-residual-transport}
For $1\leq M\leq N$, put
\[
  T_{M,N,s}(f)=\sum_{n=M+1}^{N}f(n)n^{-s}.
\]
Then
\begin{align}
\Boundary_N(f,g;s)
={}&\Boundary_M(\rho_{M,N}f,\rho_{M,N}g;s)
 +D_{M,s}(\rho_{M,N}f)T_{M,N,s}(g) \notag\\
&+D_{M,s}(\rho_{M,N}g)T_{M,N,s}(f)
 +T_{M,N,s}(f)T_{M,N,s}(g) \notag\\
&-T_{M,N,s}(f*_Ng).
\label{eq:boundary-residual-transport}
\end{align}
\end{proposition}

\begin{proof}
Write $D_{N,s}=D_{M,s}\rho_{M,N}+T_{M,N,s}$ on both scalar factors in
\cref{eq:exact-scalarisation-boundary}, subtract the corresponding identity at
resolution $M$, and use the restriction homomorphism of
\cref{prop:truncated-algebra-restriction}.
\end{proof}

\begin{proposition}[Finite additive--multiplicative ledger]
\label{prop:finite-additive-multiplicative-ledger}
Let $p_1<\cdots<p_r$ be the primes not exceeding both $y$ and $N$, and set
\[
  q_j(s)=\sum_{k=0}^{K_{p_j}(N)}p_j^{-ks}.
\]
Define the finite exponent-box overshoot
\[
  \Omega^\Box_{N,y}(s)
  =
  \sum_{\substack{0\leq k_j\leq K_{p_j}(N)\\
                  \prod_jp_j^{k_j}>N}}
  \left(\prod_jp_j^{k_j}\right)^{-s}
\]
and the unresolved-prime ancestry
\[
  U_{N,y}(s)
  =
  \sum_{\substack{n\leq N\\P^+(n)>y}}n^{-s}.
\]
Then, for every $s\in\C$,
\begin{equation}
\label{eq:finite-additive-multiplicative-ledger}
  \sum_{n\leq N}n^{-s}
  =
  \prod_{j=1}^{r}q_j(s)
  -\Omega^\Box_{N,y}(s)
  +U_{N,y}(s).
\end{equation}
\end{proposition}

\begin{proof}
Expanding the finite product gives one term for each exponent vector in the
displayed box. Unique factorisation identifies vectors whose product is at most
$N$ with the $y$-smooth integers at most $N$; all other vectors give
$\Omega^\Box_{N,y}$. The integers at most $N$ then split disjointly into the
$y$-smooth and unresolved-prime classes.
\end{proof}

\begin{proposition}[Ordinary Euler-product ledger]
\label{prop:ordinary-euler-product-ledger}
For $\Re(s)>0$, define
\[
  \Pi_y(s)=\prod_{p\leq y}(1-p^{-s})^{-1},
  \qquad
  O_{N,y}(s)=
  \sum_{\substack{n>N\\P^+(n)\leq y}}n^{-s}.
\]
Then
\begin{equation}
\label{eq:ordinary-euler-product-ledger}
  \sum_{n\leq N}n^{-s}
  =\Pi_y(s)-O_{N,y}(s)+U_{N,y}(s).
\end{equation}
The convergence boundary belongs to this scalar projection; it does not alter
the finite ancestry or \cref{eq:finite-euler-inversion}.
\end{proposition}

\begin{proof}
There are finitely many prime factors in $\Pi_y$, and each geometric series is
absolutely convergent when $\Re(s)>0$. Unique factorisation therefore gives
\[
  \Pi_y(s)=\sum_{n\in\mathcal S_y}n^{-s}.
\]
Splitting at $N$ gives the resolved $y$-smooth sum plus $O_{N,y}$; adjoining
$U_{N,y}$ recovers all integers through $N$.
\end{proof}

\begin{proposition}[Euler cutoff transport]
\label{prop:euler-cutoff-transport}
Let $m=N+1$ and
$\chi_{m,y}=\mathbf1_{\{P^+(m)\leq y\}}$. If
\[
  A_{N,y}(s)=\sum_{\substack{n\leq N\\P^+(n)\leq y}}n^{-s},
\]
then
\begin{align}
  A_{N+1,y}&=A_{N,y}+\chi_{m,y}m^{-s},\notag\\
  U_{N+1,y}&=U_{N,y}+(1-\chi_{m,y})m^{-s},\notag\\
  O_{N+1,y}&=O_{N,y}-\chi_{m,y}m^{-s}.
\label{eq:integer-cutoff-transport}
\end{align}
For $z\geq y$, put
\[
  \Delta_{N;y,z}(s)
  =\sum_{\substack{n\leq N\\y<P^+(n)\leq z}}n^{-s}.
\]
Then
\begin{align}
  A_{N,z}&=A_{N,y}+\Delta_{N;y,z},\notag\\
  U_{N,z}&=U_{N,y}-\Delta_{N;y,z},\notag\\
  O_{N,z}&=O_{N,y}+\Pi_z-\Pi_y-\Delta_{N;y,z}.
\label{eq:prime-cutoff-transport}
\end{align}
Thus changing either finite resolution transports ancestry between declared
coordinates while preserving the complete ledger.
\end{proposition}

\begin{proof}
The integer $m$ is either $y$-smooth or unresolved, giving
\cref{eq:integer-cutoff-transport}. Increasing the prime cutoff moves exactly
the integers counted by $\Delta_{N;y,z}$ from $U$ to $A$. The ordinary-product
identity \cref{eq:ordinary-euler-product-ledger} then gives the transport of
$O$.
\end{proof}

\subsection{Finite continuation transport}

Let $(s)_j=s(s+1)\cdots(s+j-1)$ and let
\(\widetilde B_{2m}(x)=B_{2m}(x-\lfloor x\rfloor)\) be the periodic Bernoulli
polynomial.

\begin{definition}[Finite continuation coordinate]
\label{def:finite-continuation-coordinate}
For $m,N\geq1$ and $s\neq1$, define
\begin{align}
  \Continuation_{m,N}(s)
  :={}&
  \sum_{n=1}^{N}n^{-s}
  +\frac{N^{1-s}}{s-1}
  -\frac12N^{-s}\notag\\
  &+
  \sum_{r=1}^{m}
  \frac{B_{2r}}{(2r)!}(s)_{2r-1}N^{1-s-2r}.
\label{eq:finite-continuation-coordinate}
\end{align}
For $1\leq N<M$, define the enacted residual
\begin{equation}
\label{eq:finite-continuation-residual}
  \mathcal U_{m;N\to M}(s)
  =
  -\frac{(s)_{2m}}{(2m)!}
  \int_N^M
  \widetilde B_{2m}(x)x^{-s-2m}\,dx.
\end{equation}
\end{definition}

\begin{proposition}[Exact finite continuation cocycle]
\label{prop:finite-continuation-cocycle}
For every finite $1\leq N<M$,
\begin{equation}
\label{eq:finite-continuation-cocycle}
  \Continuation_{m,M}(s)-\Continuation_{m,N}(s)
  =\mathcal U_{m;N\to M}(s).
\end{equation}
The residual composes exactly:
\[
  \mathcal U_{m;N\to L}
  =
  \mathcal U_{m;N\to M}
  +\mathcal U_{m;M\to L}.
\]
The finite transport is entire in $s$ after interpreting the elementary
quotient at $s=1$ by its removable value; $\Continuation_{m,N}$ is meromorphic
with its single possible pole at $s=1$.
\end{proposition}

\begin{proof}
Euler--Maclaurin on the finite half-open interval $(N,M]$ gives
\begin{align*}
\sum_{N<n\leq M}n^{-s}
={}&
\frac{M^{1-s}-N^{1-s}}{1-s}
+\frac{M^{-s}-N^{-s}}2\\
&+\sum_{r=1}^{m}
\frac{B_{2r}}{(2r)!}(s)_{2r-1}
\left(N^{1-s-2r}-M^{1-s-2r}\right)
+\mathcal U_{m;N\to M}(s).
\end{align*}
Rearrangement is \cref{eq:finite-continuation-cocycle}. Additivity of the
finite integral gives the cocycle law. No infinite endpoint is used.
\end{proof}

\begin{proposition}[Continuation as a projection of finite transport]
\label{prop:continuation-projection}
On every compact subset of $\Re(s)>1-2m$ avoiding $s=1$, the residual family
is Cauchy under sequential refinement. Define
\[
  \mathcal U_{m,N}(s)
  =
  \SeqProj_{M}
  \mathcal U_{m;N\to M}(s).
\]
Then
\begin{equation}
\label{eq:continuation-ledger}
  \zeta_{\mathrm{AC}}(s)
  =
  \Continuation_{m,N}(s)+\mathcal U_{m,N}(s),
\end{equation}
and
\[
  \mathcal U_{m,N}
  =
  \mathcal U_{m;N\to M}+\mathcal U_{m,M}
\]
for every finite $M>N$. Substituting
\cref{eq:finite-additive-multiplicative-ledger} into
\cref{eq:continuation-ledger} retains the complete arithmetic ancestry:
\begin{align}
\zeta_{\mathrm{AC}}(s)
={}&
\prod_{p_j\leq\min\{y,N\}}
\sum_{k=0}^{K_{p_j}(N)}p_j^{-ks}
-\Omega^\Box_{N,y}(s)+U_{N,y}(s)\notag\\
&+\frac{N^{1-s}}{s-1}-\frac12N^{-s}
+\sum_{r=1}^{m}
\frac{B_{2r}}{(2r)!}(s)_{2r-1}N^{1-s-2r}
+\mathcal U_{m,N}(s).
\label{eq:ancestry-bearing-continuation}
\end{align}
\end{proposition}

\begin{proof}
The periodic Bernoulli polynomial is bounded, so on a compact set with
$\Re(s)>1-2m$ the tail of
\cref{eq:finite-continuation-residual} is bounded uniformly by a constant
times
\[
  \int_N^\infty x^{-\Re(s)-2m}\,dx.
\]
This proves sequential Cauchy refinement. Where $\Re(s)>1$, the correction
terms and residual tail tend to zero while the finite sum tends to the original
Dirichlet projection. Identity of analytic continuations then gives
\cref{eq:continuation-ledger} throughout the stated domain. The cocycle retains
the unresolved continuation at every finite resolution, and substitution gives
\cref{eq:ancestry-bearing-continuation}.
\end{proof}

\begin{reviewbox}
\textbf{The static bound at $\Re(s)=1$.}
The line is the absolute-convergence boundary of the original unbounded
Dirichlet and Euler-product projections. It is not a boundary of finite natural
generation, unique factorisation, truncated convolution, Möbius inversion, or
finite continuation transport. Analytic continuation changes the projection
and carries its residual; it does not terminate or replace multiplicative
ancestry.
\end{reviewbox}

\subsection{Finite reciprocal-scale transport}

\begin{definition}[Finite theta defect]
\label{def:finite-theta-defect}
For $x>0$ and $N,M\geq0$, put
\[
  \Theta_N(x)=\sum_{n=-N}^{N}e^{-\pi n^2x},
  \qquad
  \Delta_{N,M}^{\Theta}(x)
  =
  \Theta_N(x)-x^{-1/2}\Theta_M(x^{-1}).
\]
\end{definition}

\begin{proposition}[Theta residual transport]
\label{prop:theta-residual-transport}
At every finite resolution,
\begin{equation}
\label{eq:theta-residual-transport}
  \Theta_N(x)
  =
  x^{-1/2}\Theta_M(x^{-1})
  +\Delta_{N,M}^{\Theta}(x),
\end{equation}
with
\[
  \Delta_{N,M}^{\Theta}(x)
  =
  -x^{-1/2}\Delta_{M,N}^{\Theta}(x^{-1}).
\]
If $N'\geq N$ and $M'\geq M$, then
\begin{align*}
\Delta_{N',M'}^{\Theta}(x)-\Delta_{N,M}^{\Theta}(x)
={}&2\sum_{n=N+1}^{N'}e^{-\pi n^2x}\\
&-2x^{-1/2}\sum_{k=M+1}^{M'}e^{-\pi k^2/x}.
\end{align*}
\end{proposition}

\begin{proof}
The first identity declares the exact finite residual. Replacing $x$ by
$x^{-1}$ and exchanging $N,M$ gives its involutive transport. Subtracting two
finite defects gives the refinement formula.
\end{proof}

\begin{proposition}[Finite Gaussian--Poisson bridge]
\label{prop:finite-gaussian-poisson-bridge}
Let
\[
  g_x(u)=e^{-\pi xu^2},
  \qquad
  P_{x,L}(u)=\sum_{\ell=-L}^{L}g_x(u+\ell)
\]
on $-\tfrac12\leq u\leq\tfrac12$, and put
\[
  c_{k,L}(x)
  =
  \int_{-L-1/2}^{L+1/2}
  g_x(v)e^{-2\pi ikv}\,dv.
\]
If
\[
  F_{K,L}(x)=P_{x,L}(0)-\sum_{|k|\leq K}c_{k,L}(x),
\]
then the finite theta defect satisfies
\begin{equation}
\label{eq:finite-gaussian-poisson-bridge}
  \Delta_{L,K}^{\Theta}(x)
  =
  F_{K,L}(x)
  +\sum_{|k|\leq K}
   \left(
     c_{k,L}(x)-x^{-1/2}e^{-\pi k^2/x}
   \right).
\end{equation}
Sequential enlargement of the finite intervals recovers
\begin{equation}
\label{eq:jacobi-inversion-recovered}
  \Theta(x)=x^{-1/2}\Theta(x^{-1})
\end{equation}
as the static projection of the finite residual transport.
\end{proposition}

\begin{proof}
Since $P_{x,L}(0)=\Theta_L(x)$, subtracting
\[
  x^{-1/2}\Theta_K(x^{-1})
  =\sum_{|k|\leq K}x^{-1/2}e^{-\pi k^2/x}
\]
from the definition of $F_{K,L}$ gives
\cref{eq:finite-gaussian-poisson-bridge}. The full Gaussian transform is
\[
  \int_{\mathbb R}g_x(v)e^{-2\pi ikv}\,dv
  =x^{-1/2}e^{-\pi k^2/x}.
\]
For fixed $K$ and $|k|\leq K$, the finite-interval discrepancy obeys
\[
  \left|
    c_{k,L}(x)-x^{-1/2}e^{-\pi k^2/x}
  \right|
  \leq
  \frac{e^{-\pi x(L+1/2)^2}}{\pi x(L+1/2)},
\]
so its finite sum tends to zero as $L$ is enlarged. The periodisations
$P_{x,L}$ and all their derivatives converge uniformly to the smooth
one-periodic function
\[
  P_x(u)=\SeqProj_L\sum_{\ell=-L}^{L}g_x(u+\ell).
\]
Its $k$th Fourier coefficient is the displayed full Gaussian transform.
Uniform convergence of the Fourier series of this smooth periodic function
therefore gives
\[
  \lim_{K\to\infty}
  \left(
    P_x(0)-\sum_{|k|\leq K}x^{-1/2}e^{-\pi k^2/x}
  \right)=0.
\]
Taking the declared iterated refinement $L\to\infty$ for each finite $K$ and
then $K\to\infty$ makes both terms on the right of
\cref{eq:finite-gaussian-poisson-bridge} vanish. This proves
\cref{eq:jacobi-inversion-recovered}, not merely convergence of an unspecified
residual.
\end{proof}

\begin{proposition}[Mellin recovery of functional symmetry]
\label{prop:mellin-functional-symmetry}
Let
\[
  \psi_N(x)=\frac{\Theta_N(x)-1}{2}.
\]
For finite $N,M,T>1$, define
\[
  \mathcal R_{N,M;T}(s)
  =
  \frac12\int_{1/T}^{1}
  \Delta_{N,M}^{\Theta}(x)x^{s/2-1}\,dx
\]
and
\[
  E_T(s)
  =
  \frac{1-T^{(1-s)/2}}{s-1}
  -\frac{1-T^{-s/2}}s,
\]
with removable values at $s=0,1$. Exact finite scale transport gives
\begin{align}
\int_{1/T}^{1}\psi_N(x)x^{s/2-1}\,dx
={}&
\int_1^T\psi_M(u)u^{(1-s)/2-1}\,du\notag\\
&+E_T(s)+\mathcal R_{N,M;T}(s).
\label{eq:finite-mellin-scale-transport}
\end{align}
For $M=N$, put
\[
  \mathcal M_{N,T}(s)
  =
  \int_{1/T}^{T}\psi_N(x)x^{s/2-1}\,dx
\]
and
\[
  \Boundary_T^{\mathrm M}(s)
  =
  \frac{T^{(1-s)/2}}{s-1}
  -\frac{T^{-s/2}}s.
\]
Define the finite symmetric coordinate by its transported residual ledger:
\begin{align}
  \Lambda^\sharp_{N,T}(s)
  :={}&
  \mathcal M_{N,T}(s)
  -\mathcal R_{N,N;T}(s)
  +\Boundary_T^{\mathrm M}(s)\notag\\
  ={}&
  \frac1{s-1}-\frac1s\notag\\
  &+\int_1^T\psi_N(x)
  \left(x^{s/2-1}+x^{(1-s)/2-1}\right)\,dx.
\label{eq:finite-mellin-symmetric-coordinate}
\end{align}
Then, exactly,
\[
  \Lambda^\sharp_{N,T}(s)
  =\Lambda^\sharp_{N,T}(1-s).
\]
Sequential projection of the finite theta and scale residuals gives
\begin{equation}
\label{eq:lambda-recovered}
  \Lambda(s)
  =
  \frac1{s(s-1)}
  +\SeqProj_T
  \int_1^T\psi(x)
  \left(x^{s/2-1}+x^{(1-s)/2-1}\right)\,dx,
\end{equation}
where
\[
  \Lambda(s)=\pi^{-s/2}\Gamma(s/2)\zeta_{\mathrm{AC}}(s).
\]
Consequently
\begin{equation}
\label{eq:functional-symmetry-recovered}
  \xi(s)=\xi(1-s).
\end{equation}
\end{proposition}

\begin{proof}
Rewrite \cref{eq:theta-residual-transport} as
\[
  \psi_N(x)
  =
  x^{-1/2}\psi_M(x^{-1})
  +\frac{x^{-1/2}-1}{2}
  +\frac{\Delta_{N,M}^{\Theta}(x)}2.
\]
Multiplication by $x^{s/2-1}$, integration on $[1/T,1]$, and the substitution
$u=x^{-1}$ give \cref{eq:finite-mellin-scale-transport}; direct integration of
the zero-mode term gives $E_T(s)$. Taking $M=N$ and adjoining the upper-scale
integral proves both expressions for $\Lambda^\sharp_{N,T}$. The two integrand
terms exchange under $s\mapsto1-s$, as do $1/(s-1)$ and $-1/s$. This proves
finite symmetry before any limit.

For $\Re(s)>1$, termwise Mellin integration of the finite Gaussian sums and
the residual transport \cref{eq:theta-residual-transport} give
\[
  \pi^{-s/2}\Gamma(s/2)\zeta(s)
  =\int_0^\infty\psi(x)x^{s/2-1}\,dx.
\]
Splitting at $1$, applying recovered Jacobi inversion below $1$, and changing
$x\mapsto x^{-1}$ gives \cref{eq:lambda-recovered}; the removed zero mode
contributes $1/(s(s-1))$. The expression provides continuation and is invariant
under $s\mapsto1-s$. Multiplication by
$\tfrac12s(s-1)$ proves \cref{eq:functional-symmetry-recovered}.
\end{proof}

Theta, Mellin, Fourier, heat-flow, product, and prime-side representations are
projection coordinates of the same relation. Their roles follow from exact
dependency rather than classical naming, symmetry, or convenience.

\section{The invariant process}
\label{sec:process-invariant}

\begin{definition}[Finite Möbius--hyperbola invariant]
\label{def:finite-moebius-hyperbola}
For real $X\geq1$ and $s\in\C$, define
\begin{equation}
\label{eq:finite-moebius-hyperbola}
  \Invariant_X(s)
  =
  \sum_{\substack{a,b\in\N\\ab\leq X}}
  \frac{\mu(a)}{(ab)^s}.
\end{equation}
\end{definition}

\begin{proposition}[Finite invariant normalisation]
\label{prop:finite-moebius-hyperbola}
For every $X\geq1$ and every $s\in\C$,
\begin{equation}
\label{eq:finite-invariant-normalisation}
  \Invariant_X(s)=1.
\end{equation}
The identity is finite, entire in $s$, and independent of numerical magnitude.
\end{proposition}

\begin{proof}
Regrouping the finite sum by $n=ab$ gives
\[
  \Invariant_X(s)
  =
  \sum_{n\leq X}n^{-s}\sum_{a\mid n}\mu(a).
\]
The inner sum is $1$ for $n=1$ and $0$ otherwise by
\cref{eq:finite-moebius-inverse}. Since $X\geq1$, the result is $1$.
\end{proof}

For $y\geq0$, set
\[
  Z_s(y)=\sum_{n\leq y}n^{-s},
  \qquad
  M_s(y)=\sum_{n\leq y}\mu(n)n^{-s},
\]
with either sum zero when $y<1$.

\begin{proposition}[Arbitrary hyperbola cutoff]
\label{prop:arbitrary-hyperbola-cutoff}
Let $X\geq1$ and let $u,v>0$ satisfy $uv=X$. Then
\begin{align}
  \Invariant_X(s)
  ={}&
  \sum_{a\leq u}\frac{\mu(a)}{a^s}
    Z_s\!\left(\frac Xa\right)
  +\sum_{b\leq v}\frac1{b^s}
    M_s\!\left(\frac Xb\right)\notag\\
  &-M_s(u)Z_s(v)
  =1.
\label{eq:arbitrary-hyperbola-cutoff}
\end{align}
\end{proposition}

\begin{proof}
Every point $ab\leq X$ lies in $a\leq u$ or $b\leq v$, since otherwise
$ab>uv=X$. Their intersection is the complete rectangle
$a\leq u$, $b\leq v$. Inclusion--exclusion gives the three terms in
\cref{eq:arbitrary-hyperbola-cutoff}; their value is
\cref{eq:finite-invariant-normalisation}.
\end{proof}

\begin{corollary}[Exact hyperbola transport]
\label{cor:exact-hyperbola-transport}
For two cutoff pairs $u_jv_j=X$, let $A(u)$, $B(v)$, and $C(u,v)$ denote the
three respective terms of \cref{eq:arbitrary-hyperbola-cutoff}. Then
\[
  [A(u_1)-A(u_0)]
  +[B(v_1)-B(v_0)]
  -[C(u_1,v_1)-C(u_0,v_0)]
  =0.
\]
Changing resolution transports the complete boundary ledger without a
discarded error.
\end{corollary}

\begin{proposition}[Unique self-dual resolution]
\label{prop:unique-self-dual-cutoff}
Fix $X>1$ and parameterise the positive cutoff pairs by
\[
  u_\theta=X^\theta,
  \qquad
  v_\theta=X^{1-\theta},
  \qquad
  \theta\in\mathbb R.
\]
Exchange of the factor channels sends $\theta\mapsto1-\theta$. Its unique fixed
resolution is
\begin{equation}
\label{eq:unique-self-dual-cutoff}
  \theta=\frac12,
  \qquad
  u_\theta=v_\theta=\sqrt X.
\end{equation}
The resolution is independent of $X$, $s$, and the attained numerical
magnitude.
\end{proposition}

\begin{proof}
The cutoff pair is fixed precisely when
$X^\theta=X^{1-\theta}$. Since $X>1$, taking logarithms gives
$\theta=1-\theta$, and hence $\theta=\tfrac12$.
\end{proof}

\begin{remark}[Continuous and integer resolutions]
\label{rem:cutoff-uniqueness-scope}
The uniqueness is in the positive real cutoff coordinate. At $X=1$, every
exponent gives the same unit cutoff. After floors, distinct exponents may also
name the same finite integer pair; an exactly equal integer pair exists only
when $X$ is a square. These finite aliases do not alter the unique continuous
fixed resolution.
\end{remark}

\begin{proposition}[Inverse-convolution widening]
\label{prop:inverse-convolution-universality}
Let arithmetic functions $f,g$ satisfy $f*g=\varepsilon$. Then
\[
  \Invariant_X^{f,g}(s)
  =
  \sum_{ab\leq X}\frac{f(a)g(b)}{(ab)^s}
  =1
\]
for every $X\geq1$ and $s\in\C$.
\end{proposition}

\begin{proof}
Regrouping by $n=ab$ gives
\[
  \Invariant_X^{f,g}(s)
  =\sum_{n\leq X}\frac{(f*g)(n)}{n^s}
  =\sum_{n\leq X}\frac{\varepsilon(n)}{n^s}
  =1.
\]
\end{proof}

\begin{keybox}[title={\bfseries Process invariant and its projection}]
The invariant is the complete additive--multiplicative boundary ledger
\cref{eq:arbitrary-hyperbola-cutoff}, not either separated channel. Its factor
exchange has the magnitude-independent resolution
\cref{eq:unique-self-dual-cutoff}. Jacobi--Mellin transport independently maps
reciprocal scale into $s\leftrightarrow1-s$ by
\cref{prop:mellin-functional-symmetry}. The functional symmetry is therefore a
static projection of self-dual arithmetic and scale transport rather than the
primitive relation.
\end{keybox}

\section{Counterexample differential and carrier exhaustion}
\label{sec:counterexample-differential}

\begin{definition}[Off-line-zero differential]
\begin{equation}
\label{eq:off-line-zero-differential}
  \CounterDiff:
  \qquad
  \xi(\rho)=0,\quad0<\Re(\rho)<1,\quad\rho\neq\Dual(\rho).
\end{equation}
This is a hypothetical finite witness, not a completed infinite branch.
\end{definition}

\begin{definition}[Twisted Möbius coordinate]
\label{def:twisted-mobius-coordinate}
For $\gamma\in\mathbb R$ and $X\geq1$, set
\[
  M_\gamma(X)=\sum_{n\leq X}\mu(n)n^{-i\gamma}.
\]
The untwisted coordinate $M_0(X)$ is the Mertens function $M(X)$.
\end{definition}

\begin{proposition}[Finite Möbius--Mellin ledger]
\label{prop:finite-mobius-mellin-ledger}
For $X\geq1$ and $z\in\C$,
\begin{equation}
\label{eq:finite-mobius-mellin-ledger}
  \sum_{n\leq X}\frac{\mu(n)}{n^{z+i\gamma}}
  =
  M_\gamma(X)X^{-z}
  +z\int_1^X M_\gamma(x)x^{-z-1}\,dx.
\end{equation}
For $\Re(z)>1$, sequential projection gives
\[
  \frac1{\zeta(z+i\gamma)}
  =
  z\int_1^\infty M_\gamma(x)x^{-z-1}\,dx.
\]
In the same domain, at every finite $X$, the exact unresolved boundary is
\begin{equation}
\label{eq:mobius-mellin-residual}
  \Residual_{\gamma,z}(X)
  =
  -M_\gamma(X)X^{-z}
  +z\int_X^\infty M_\gamma(x)x^{-z-1}\,dx.
\end{equation}
\end{proposition}

\begin{proof}
Apply finite Abel summation to
$a_n=\mu(n)n^{-i\gamma}$ and $f(x)=x^{-z}$:
\[
  \sum_{n\leq X}a_nf(n)
  =M_\gamma(X)f(X)-\int_1^XM_\gamma(x)f'(x)\,dx.
\]
This is \cref{eq:finite-mobius-mellin-ledger}. In $\Re(z)>1$, absolute
convergence identifies the sequential scalar projection with
$1/\zeta(z+i\gamma)$. Splitting the projected integral at finite $X$ gives
\cref{eq:mobius-mellin-residual}; no finite boundary term is deleted.
\end{proof}

\begin{theorem}[Möbius growth obstruction from a zero]
\label{thm:mobius-growth-obstruction}
If $\zeta(\beta+i\gamma)=0$, then for every $\varepsilon>0$,
\begin{equation}
\label{eq:twisted-growth-obstruction}
  M_\gamma(X)\neq O\!\left(X^{\beta-\varepsilon}\right).
\end{equation}
The untwisted Mertens function also satisfies
\[
  M(X)\neq O\!\left(X^{\beta-\varepsilon}\right).
\]
\end{theorem}

\begin{proof}
If $M_\gamma(X)=O(X^{\beta-\varepsilon})$, then
\[
  F(z)=z\int_1^\infty M_\gamma(x)x^{-z-1}\,dx
\]
converges locally uniformly and is holomorphic for
$\Re(z)>\beta-\varepsilon$. In $\Re(z)>1$ it equals
$1/\zeta(z+i\gamma)$ by
\cref{prop:finite-mobius-mellin-ledger}. Uniqueness of analytic continuation
would make the reciprocal holomorphic at $z=\beta$, contradicting the zero at
$\beta+i\gamma$. The untwisted proof applies the same argument to
\[
  \frac1{\zeta(s)}
  =s\int_1^\infty M(x)x^{-s-1}\,dx.
\]
\end{proof}

\begin{definition}[Persistent exponent differential]
\label{def:persistent-exponent-differential}
Let an off-line zero be represented, using conjugate reflection, as
\[
  \rho=\frac12+\delta+i\gamma,
  \qquad
  \delta>0.
\]
Its persistent exponent differential is
\begin{equation}
\label{eq:persistent-exponent-differential}
  \mathcal P_\rho:
  \qquad
  \forall\,0<\varepsilon<\delta,\quad
  \limsup_{X\to\infty}
  \frac{|M_\gamma(X)|}
       {X^{1/2+\delta-\varepsilon}}
  =\infty.
\end{equation}
It records the obstruction $X^\delta$ relative to the self-dual
$X^{1/2}$ resolution; it does not assert a pointwise sign or an asymptotic
formula for $M_\gamma$.
\end{definition}

\begin{corollary}[Off-line zero implies exponent persistence]
\label{cor:off-line-zero-implies-persistence}
For every off-line nontrivial zero $\rho$, there is a right-hand representative
\[
  \rho^+\in\{\rho,\Dual(\rho)\},
  \qquad
  \rho^+=\frac12+\delta+i\gamma,
  \quad\delta>0,
\]
and $\rho^+$ satisfies $\mathcal P_{\rho^+}$.
\end{corollary}

\begin{proof}
The symmetry recovered in \cref{prop:mellin-functional-symmetry}, together with
real conjugation, gives the reflected zero orbit. Choose $\rho^+$ to the right
of $\CriticalLine$. Then
\cref{thm:mobius-growth-obstruction} with
$\beta=\tfrac12+\delta$ gives
\cref{eq:persistent-exponent-differential}.
\end{proof}

\begin{definition}[Zeta carrier audit]
\label{def:zeta-carrier-audit}
Write $\Audit_\zeta(\rho)$ when the following carriers have all been retained
with their dependency paths:
\begin{enumerate}[label=(\arabic*)]
  \item finite integers, prime ancestry, divisor pairs, $\mu$, $\mathbf1$, and
        $\varepsilon$;
  \item the hyperbola cutoffs, overlap rectangle, factor-channel exchange, and
        finite floor aliases;
  \item prime cutoff, exponent-box overshoot, unresolved-prime ancestry, and
        scalarisation boundary;
  \item $M(X)$, $M_\gamma(X)$, their Abel boundary terms, and Mellin residuals;
  \item finite theta truncation, Gaussian--Poisson residual, Mellin split, and
        polar zero-mode term;
  \item continuation coordinates, Bernoulli residuals, and the order of every
        finite refinement;
  \item completion factors, logarithm branch, and the full reflection orbit;
  \item height, phase, and zero multiplicity;
  \item smoothing weights, kernels, contours, and regularisations, recorded as
        absent from this derivation rather than silently assigned zero
        residuals; and
  \item the spectral--hyperbola projection, self-dual cutoff and Mellin
        projections, raw finite registration family, generated context and
        resolution, coherence and exponent-coordinate projections, constructor
        algebra, finite branch family, threshold, persistence weights, and
        complete normalisation record.
\end{enumerate}
\end{definition}

\begin{proposition}[Carrier ancestry closure]
\label{prop:zeta-carrier-ancestry}
For every off-line-zero differential
\cref{eq:off-line-zero-differential}, every carrier in
$\Audit_\zeta(\rho)$ retains ancestry to the finite generator
\cref{def:finite-prime-ancestry,def:truncated-dirichlet-algebra} or to an exact
projection and residual transport derived from it. Each carrier used in the
proof is an explicit deterministic function of registered variables; the
candidate carriers not used in the proof are explicitly absent.
\end{proposition}

\begin{proof}
Items (1)--(3) are finite consequences of unique factorisation,
\cref{prop:finite-euler-inversion,prop:exact-scalarisation-boundary,prop:euler-cutoff-transport}.
Item (4) is the finite Abel ledger
\cref{prop:finite-mobius-mellin-ledger}. Items (5)--(6) are the exact theta,
Mellin, and continuation transports
\cref{prop:finite-continuation-cocycle,prop:theta-residual-transport,prop:finite-gaussian-poisson-bridge}.
The completion factors in item (7) are explicit and nonzero inside the
nontrivial strip. The branch is fixed by
$n^{-s}=\exp(-s\log n)$ with the real logarithm of $n>0$, and the reflection
orbit is generated by \cref{eq:functional-symmetry-recovered} and complex
conjugation. In item (8), height is the parameter $\gamma$ and phase is the
explicit character
\[
  n^{-i\gamma}=e^{-i\gamma\log n}.
\]
A zero of any positive multiplicity makes $1/\zeta$ non-holomorphic, so
\cref{thm:mobius-growth-obstruction} assumes no simple zero and its exponent
obstruction is unchanged. Item (9) lists operations that do not occur in the
derivation: no smoothing kernel, contour displacement, or regularisation is
used to obtain \cref{thm:mobius-growth-obstruction}. They therefore contribute
no undeclared terms. Item (10) is constructed in
\cref{def:riemann-process-field,def:finite-factor-exchange-branch,def:riemann-persistence-coordinates,def:riemann-registration-family,def:finite-riemann-persistence,def:per-zero-normalisation-record}
from the spectral coordinate, exact finite ledgers, and the single Cogenesis
persistence relation; each projection states its residual and no second
normaliser is introduced. This proves the dependency registry only; its
extension through causal completion occurs solely in
\cref{rev:foundational-inference}.
\end{proof}

\begin{reviewbox}
\textbf{Status of carrier exhaustion.}
The audit does not erase the off-line candidate. It retains its exact
$X^\delta$ obstruction, every residual through which it is expressed, and the
unchanged ancestry that generates those coordinates. The foundational
causal extension in \cref{rev:foundational-inference} determines whether that
fully typed differential completes as the realised output or remains residual.
\end{reviewbox}

\section{Persistence selection in the Riemann register}
\label{sec:riemann-persistence}

\subsection{Refinement-defined zero projection}

\begin{definition}[Zero-refinement projection]
\label{def:zero-refinement-projection}
For a hypothetical zero
$\rho=\tfrac12+\delta+i\gamma$ and each finite refinement record $\mathbf f$,
put $z_\rho=\tfrac12+\delta$ and define
\[
  \mathcal O_{\Dual}(\rho)
  :=
  \{\rho,\overline\rho,1-\rho,1-\overline\rho\},
  \qquad
  \mathfrak A_{\gamma,z_\rho;X}
  :=
  \left(
    M_\gamma(X)X^{-z_\rho},
    z_\rho\int_1^X M_\gamma(x)x^{-z_\rho-1}\,dx
  \right).
\]
The exact initial typed record is
\[
  \Field_{\rho,\mathbf f}^{(0)}
  :=
  \left(
    \mathcal O_{\Dual}(\rho),
    \operatorname{Comp}(\rho);
    \mathscr L_{\mathbf f};
    M_\gamma(X),\mathfrak A_{\gamma,z_\rho;X};
    \Residual_{\rho;\mathbf f}
  \right),
\]
where $\operatorname{Comp}(\rho)$ is the explicit tuple of completion factors
in \cref{eq:completed-xi} and $\Residual_{\rho;\mathbf f}$ is the tuple of
their separately typed finite projection and continuation residuals. For
$n\in\mathbb N_0:=\{0,1,2,\ldots\}$, recursively define
\[
\begin{aligned}
  \Field_{\rho,\mathbf f}^{(n+1)}
  &:=
  \Field_{\rho,\mathbf f}^{(n)}
  +
  \left\{
    G(x)+R_G(x):
    G\in\operatorname{Op}(\Gen_{\zeta,\mathbf f}),\
    x\in
    \operatorname{Carr}(\Field_{\rho,\mathbf f}^{(n)})
    \cap\operatorname{Dom}(G)
  \right\}.
\end{aligned}
\]
where $R_G(x)$ is the declared residual of the finite operation $G$. Then
\[
  \Field_\rho
  :=
  \left(
    \{\Field_{\rho,\mathbf f}^{(n)}\}_{\mathbf f\ {\rm finite},\,n\in\mathbb N_0},
    \{\mathsf T_{\mathbf f\to\mathbf f'}\}
  \right).
\]
This is the smallest sequential $\Gen_\zeta$-generated projection: every
finite generator word occurs at a finite $n$, while no completed union over
$n$ or $\mathbf f$ is formed. Indeed, any domain-respecting
$\Gen_{\zeta,\mathbf f}$-closed record containing
$\Field_{\rho,\mathbf f}^{(0)}$ contains every finite stage by induction.
For a finite $R\geq1$, the tail projection $\tau_R$ acts on the twisted
Möbius coordinate by
\[
  \tau_RM_\gamma(X)
  =
  M_\gamma(X)-M_\gamma(R),
  \qquad X\geq R,
\]
and transports every other coordinate by its already derived finite cocycle.
\end{definition}

\begin{lemma}[Finite-prefix invariance of exponent persistence]
\label{lem:riemann-tail-stability}
If $\mathcal P_\rho$ of
\cref{eq:persistent-exponent-differential} holds, then for every finite
$R\geq1$ and every $0<\varepsilon<\delta$,
\begin{equation}
\label{eq:riemann-tail-stability}
  \limsup_{X\to\infty}
  \frac{|M_\gamma(X)-M_\gamma(R)|}
       {X^{1/2+\delta-\varepsilon}}
  =\infty.
\end{equation}
Thus the off-line exponent obstruction is retained by every finite refinement
tail of $\Field_\rho$.
\end{lemma}

\begin{proof}
If the quotient in \cref{eq:riemann-tail-stability} were bounded, then adding
the fixed constant $M_\gamma(R)$ would give
\[
  M_\gamma(X)=O\!\left(X^{1/2+\delta-\varepsilon}\right),
\]
contradicting $\mathcal P_\rho$. Removing any other finite initial ledger
changes only finitely many enacted coordinates and its exact transported
boundary, so it cannot alter the asymptotic exponent statement.
\end{proof}

\begin{proposition}[Cogenetic typing of the exponent differential]
\label{prop:riemann-differential-typing}
Within the Cogenesis causal grammar
\cref{eq:imported-persistence,eq:imported-branch-normalisation},
$\mathcal P_\rho$ is:
\begin{enumerate}[label=(\roman*)]
  \item a projected expression with its Möbius--Mellin residual;
  \item a persistent bifurcating differential, because every finite tail
        retains the same strict exponent distinction; and
  \item a candidate unary of $\Field_\rho$ to be evaluated by the one
        normalising relation $\heartsim$.
\end{enumerate}
Candidate status does not identify the differential with the selected unary.
\end{proposition}

\begin{proof}
The projected expression and residual are
\cref{eq:finite-mobius-mellin-ledger,eq:mobius-mellin-residual}.
\Cref{lem:riemann-tail-stability} proves persistence under every finite tail.
Its nonzero exponent magnitude and ancestry-derived alternatives are enacted
explicitly at \cref{eq:riemann-branch-enactment}. The Cogenesis finite
persistence relation therefore types the differential as an admissible candidate while
\cref{eq:imported-branch-normalisation} retains selected, unselected, and
continuing content separately.
\end{proof}

\subsection{Exact scale residual and finite persistence}

\begin{definition}[Finite factor-exchange branch]
\label{def:finite-factor-exchange-branch}
Fix $X,T>1$ and $\gamma\in\mathbb R$. Before choosing an admissible branch,
define the raw exchange coordinate for every
\[
  \eta\in I_{\mathrm{strip}}:=\left(-\frac12,\frac12\right)
\]
put
\[
  s_\eta=\frac12+\eta+i\gamma,
\]
Declare the exponent-preserving spectral--hyperbola projection
\begin{equation}
\label{eq:spectral-hyperbola-projection}
  \Pi_X^{\mathrm{HM}}(s)
  =
  \left(X^{\Re(s)},X^{1-\Re(s)}\right),
  \qquad
  (u_\eta,v_\eta)
  :=
  \Pi_X^{\mathrm{HM}}(s_\eta).
\end{equation}
This is the domain projection carrying the real Mellin exponent into a
multiplicative cutoff pair; it is declared for every $s$ before a zero
condition is imposed.
Then $u_\eta v_\eta=X$. The exchange involution is
\[
  \tau(\eta)=-\eta,
  \qquad
  (u_\eta,v_\eta)\longmapsto(v_\eta,u_\eta),
  \qquad
  \Dual(s_\eta)=s_{-\eta}.
\]
The central record $\eta=0$ is the self-dual projection of the exchanged
cutoff pair. It is an admissible resolution record, not an assumed zero at
$\tfrac12+i\gamma$.
\end{definition}

\begin{definition}[Hyperbola coherence and scale residuals]
\label{def:riemann-persistence-coordinates}
For $uv=X$, define the complete hyperbola ledger
\begin{align*}
  \mathcal H_X(s;u,v)
  :={}&
  \sum_{a\leq u}\frac{\mu(a)}{a^s}
    Z_s\!\left(\frac Xa\right)
  +\sum_{b\leq v}\frac1{b^s}
    M_s\!\left(\frac Xb\right)\\
  &-M_s(u)Z_s(v).
\end{align*}
For the branch $\eta$, set
\[
  c_X(\eta)=\mathcal H_X(s_\eta;u_\eta,v_\eta).
\]
The self-dual cutoff projection is
\[
  \pi_{\mathrm{sd}}(u_\eta,v_\eta)
  =(\sqrt X,\sqrt X),
\]
and its retained oriented residual is
\[
  r_X^{\mathrm H}(\eta)
  =
  \frac12\log_X\!\left(\frac{u_\eta}{v_\eta}\right).
\]
For the paired Mellin kernel weights
\[
  K_T^+(s)=T^{s/2-1},
  \qquad
  K_T^-(s)=T^{(1-s)/2-1},
\]
define
\[
  \pi_{\mathrm{M,sd}}(A,B)=(\sqrt{AB},\sqrt{AB}),
  \qquad A,B>0,
\]
applied to
$A=|K_T^+(s_\eta)|$ and $B=|K_T^-(s_\eta)|$, and retain the oriented residual
\[
  r_T^{\mathrm M}(\eta)
  =
  \log_T
  \frac{|K_T^+(s_\eta)|}{|K_T^-(s_\eta)|}.
\]
\end{definition}

\begin{proposition}[Exact cross-projection exponent agreement]
\label{prop:exact-branch-agreement}
For every permitted $X,T,\gamma$ and $\eta\in I_{\mathrm{strip}}$,
\begin{equation}
\label{eq:exact-branch-agreement}
  c_X(\eta)=1,
  \qquad
  r_X^{\mathrm H}(\eta)
  =r_T^{\mathrm M}(\eta)
  =\eta.
\end{equation}
Under the declared projection \cref{eq:spectral-hyperbola-projection}, the
hyperbola cutoff coordinate and Mellin spectral coordinate carry the same
dimensionless residual without using a zero condition.
\end{proposition}

\begin{proof}
\Cref{prop:arbitrary-hyperbola-cutoff} gives $c_X(\eta)=1$ for every $s_\eta$
and every cutoff pair with product $X$. Also
\[
  \frac{u_\eta}{v_\eta}=X^{2\eta},
\]
so $r_X^{\mathrm H}(\eta)=\eta$. Finally,
\[
  |K_T^+(s_\eta)|=T^{-3/4+\eta/2},
  \qquad
  |K_T^-(s_\eta)|=T^{-3/4-\eta/2},
\]
whose ratio is $T^\eta$. Thus $r_T^{\mathrm M}(\eta)=\eta$ independently of
$X,T,\gamma$, and numerical height.
\end{proof}

\begin{definition}[Riemann registration family]
\label{def:riemann-registration-family}
Let $\mathbf f$ denote the finite arithmetic, prime, continuation, theta, and
transform cutoffs already enacted, together with finite $X,T>1$,
$1\leq R\leq X$ and height $\gamma$. Before context, resolution, or an admissible branch is
selected, form the raw finite record
\[
\begin{aligned}
  E_{\eta;\mathbf f}
  :=
  \bigl(&s_\eta,u_\eta,v_\eta,
    \mathcal H_X(s_\eta;u_\eta,v_\eta),
    K_T^+(s_\eta),K_T^-(s_\eta),
    M_\gamma(X)-M_\gamma(R);\\
  &\mathsf{Anc}_{\zeta,\mathbf f},\Residual_{\eta;\mathbf f}\bigr),
  \qquad \eta\in I_{\mathrm{strip}}.
\end{aligned}
\]
and let
$\mathscr E_{\mathbf f}
=\{E_{\eta;\mathbf f}:\eta\in I_{\mathrm{strip}}\}$.
Here $\Residual_{\eta;\mathbf f}$ is the tuple of the already-derived finite
boundary and transport residuals evaluated at $s_\eta$ or
$z_\eta:=\tfrac12+\eta$, as required by their domains. The finite family
presently capable of registering distinctions is
\[
  \mathcal G_{\mathbf f}
  =
  \left\{
    G_s,G_{\mathrm H},G_{\mathrm M},G_c,G_{\mathrm{tail},R},G_{\mathrm{meta}}
  \right\}.
\]
Their common domain is $\mathscr E_{\mathbf f}$, and their actions and
codomains are
\[
  \begin{aligned}
  G_s &: \mathscr E_{\mathbf f}\to\mathbb R,
  &G_s(E_{\eta;\mathbf f})&=\Re(s_\eta)-\frac12,\\
  G_{\mathrm H}&:\mathscr E_{\mathbf f}\to\mathbb R,
  &G_{\mathrm H}(E_{\eta;\mathbf f})&=r_X^{\mathrm H}(\eta),\\
  G_{\mathrm M}&:\mathscr E_{\mathbf f}\to\mathbb R,
  &G_{\mathrm M}(E_{\eta;\mathbf f})&=r_T^{\mathrm M}(\eta),\\
  G_c&:\mathscr E_{\mathbf f}\to\mathbb C,
  &G_c(E_{\eta;\mathbf f})&=c_X(\eta),\\
  G_{\mathrm{tail},R}&:\mathscr E_{\mathbf f}\to\mathscr T_{\mathbf f},
  &G_{\mathrm{tail},R}(E_{\eta;\mathbf f})
  &=
  \Bigl(
    M_\gamma(X)-M_\gamma(R),\,
    M_\gamma(X)X^{-z_\eta}-M_\gamma(R)R^{-z_\eta},\\
  &&&\hspace{4.4cm}
    z_\eta\int_R^X M_\gamma(x)x^{-z_\eta-1}\,dx
  \Bigr),\\
  G_{\mathrm{meta}}&:\mathscr E_{\mathbf f}\to\mathscr M_{\mathbf f},
  &G_{\mathrm{meta}}(E_{\eta;\mathbf f})
  &=
  \bigl(\mathsf{Anc}_{\zeta,\mathbf f},\Residual_{\eta;\mathbf f}\bigr),
  \end{aligned}
\]
where $\mathscr T_{\mathbf f}$ is the finite tail-ledger codomain and
$\mathscr M_{\mathbf f}$ is the finite ancestry-and-residual codomain. The
last two components of $G_{\mathrm{tail},R}$ are the exact endpoint and
integral terms obtained by subtracting
\cref{eq:finite-mobius-mellin-ledger} at $R$ from its value at $X$. These registries
are restrictions of the prior finite ancestry, not projections authorised by
the frame they generate. Following the observation-generation relation of
Cogenesis, define
\begin{equation}
\label{eq:riemann-context-resolution-generation}
  \Cctx_{\mathbf f}
  :=
  \operatorname{supp}\mathcal G_{\mathbf f}
  =\mathscr E_{\mathbf f},
  \qquad
  \Rres_{\mathbf f}
  :=
  \operatorname{Disc}\!\left(
    \mathcal G_{\mathbf f}\mid\Cctx_{\mathbf f}
  \right),
  \qquad
  \mathcal O_{\mathbf f}
  :=
  \pi_{\Cctx_{\mathbf f},\Rres_{\mathbf f}}
  (\mathcal G_{\mathbf f}).
\end{equation}
Explicitly, $\Rres_{\mathbf f}$ distinguishes
$E_{\eta;\mathbf f}$ and $E_{\eta';\mathbf f}$ whenever at least one registry
has unequal values, and
$\mathcal O_{\mathbf f}(E_{\eta;\mathbf f})$ denotes the tuple of all six
registered values.
Thus context is the support on which these finite registrations are jointly
defined, resolution is their enacted distinguishability, and observation is
their carried projection with all omitted content retained in the registered
residuals.
The exponent-channel restriction is a projection of that observation:
\[
  \pi_{\mathrm{exp}}
  \bigl(\mathcal O_{\mathbf f}(E_{\eta;\mathbf f})\bigr)
  :=
  \mathcal O_{\mathbf f}^{\mathrm{exp}}(\eta)
  +R_{\mathrm{exp};\mathbf f};
\]
$R_{\mathrm{exp};\mathbf f}$ retains $G_c$, $G_{\mathrm{tail},R}$ and
$G_{\mathrm{meta}}$, including ancestry and every non-exponent residual,
rather than deleting them.
\end{definition}

\begin{proposition}[Registration-derived exponent resolution]
\label{prop:riemann-registration-recovery}
The exponent-channel projection of
\cref{eq:riemann-context-resolution-generation} is
\begin{equation}
\label{eq:registered-exponent-observation}
  \mathcal O_{\mathbf f}^{\mathrm{exp}}(\eta)
  =
  \bigl(
    G_s(E_{\eta;\mathbf f}),
    G_{\mathrm H}(E_{\eta;\mathbf f}),
    G_{\mathrm M}(E_{\eta;\mathbf f})
  \bigr)
  =
  (\eta,\eta,\eta).
\end{equation}
Consequently exact finite registration distinguishes every nonzero $d$, gives
$\Theta_\odot=0$ in the exponent channel.
\end{proposition}

\begin{proof}
The spectral definition gives
\[
  G_s(E_{\eta;\mathbf f})=\eta,
\]
while \cref{prop:exact-branch-agreement} gives
\[
  G_{\mathrm H}(E_{\eta;\mathbf f})
  =G_{\mathrm M}(E_{\eta;\mathbf f})=\eta.
\]
Hence distinct exponent
coordinates have distinct finite registration triples.
\end{proof}

\begin{definition}[Finite Riemann persistence sample]
\label{def:finite-riemann-persistence}
For every $\eta\in I_{\mathrm{strip}}$, let $\Phi_\eta$ be the complete finite
record carrying $s_\eta$,
$(u_\eta,v_\eta)$, $\mathcal H_X$, $c_X(\eta)$, both Mellin kernels and scale
residuals, $M_\gamma(X)$, the finite Abel ledger and its boundary, and its
exchange ancestry.

For a right-hand hypothetical zero
$\rho=\tfrac12+\delta+i\gamma$, set $d=\delta$. For every finite $R\geq1$,
the paired finite projections are
\[
  q_{X,T,R}^{\pm}:\Field_\rho\longrightarrow
  \Phi_{\pm d}+R_{\rho;X,T,R}^{\pm}.
\]
The record $\Phi_d$ carries $\mathcal P_\rho$, and $\Phi_{-d}$ carries its
reflected relation. Each $R_{\rho;X,T,R}^{\pm}$ retains every component of
$\Field_\rho$ not carried by the corresponding finite record.

The lifted self-dual projection acts on the complete paired records:
\begin{equation}
\label{eq:lifted-self-dual-projection}
  \widehat\pi_{\mathrm{sd}}
  \left(
    \Phi_d+R_{\rho;X,T,R}^{+},
    \Phi_{-d}+R_{\rho;X,T,R}^{-}
  \right)
  =
  \Phi_0+\widehat R_{\rho;X,T,R}.
\end{equation}
It sends $s_{\pm d}$ to $s_0$, each cutoff pair through
$\pi_{\mathrm{sd}}$, and each Mellin magnitude pair through
$\pi_{\mathrm{M,sd}}$; $\widehat R_{\rho;X,T,R}$ retains every component not
carried by those self-dual projections.

Define the coherence and exponent-coordinate projections by
\[
  \pi_c(\Phi_\eta)=c_X(\eta),
  \qquad
  \rho_{\Phi_\eta}=|\pi_c(\Phi_\eta)|^2,
  \qquad
  \pi_\delta(\Phi_\eta)
  =r_X^{\mathrm H}(\eta)
  =r_T^{\mathrm M}(\eta)
  =\eta.
\]
By \cref{prop:riemann-registration-recovery}, every $d>0$ is distinguishable
at the registration-derived exponent resolution, so
\[
  |\DeltaFork(\Field_\rho)|:=d>\Theta_\odot:=0
\]
Every arithmetic, continuation, theta, tail, and coherence transport acts at
fixed $\eta$ in the exponent channel. The only ancestry maps that alter that
coordinate are the derived exchange and the residual-retaining self-dual
projection. Consequently the induced unary-forming constructor algebra in
$\mathsf{Anc}_{\zeta,\mathbf f}$ is
\[
  \mathfrak S_{\pitchfork}
  =
  \{\iota,\tau,\nu\},
  \qquad
  \iota(\eta)=\eta,
  \quad
  \tau(\eta)=-\eta,
  \quad
  \nu(\eta)=0.
\]
Here $\iota$ carries the source coordinate, $\tau$ is the derived exchange,
and $\nu$ is the coordinate action of
\cref{eq:lifted-self-dual-projection}. This set is closed under composition,
and therefore its orbit is derived exactly:
\[
  \mathfrak S_{\pitchfork}\!\cdot d=\{-d,0,d\}.
\]
The enacted branch relation is consequently
\begin{equation}
\label{eq:riemann-branch-enactment}
  \DeltaFork(\Field_\rho)
  \longrightarrow
  \left(
    \mathcal P_\rho;
    \{\Phi_\eta:\eta\in\mathfrak S_{\pitchfork}\!\cdot d\}
  \right).
\end{equation}
The branch context is the restriction of
$\mathcal O_{\mathbf f}^{\mathrm{exp}}$ to that orbit, with every non-exponent
registration retained in $R_{\mathrm{exp};\mathbf f}$; denote its frame by
$(\Cctx_{X,T,\gamma,d},\Rres_{\mathrm{exp}})$. Hence
\[
  \Admissible(\Cctx_{X,T,\gamma,d},\Rres_{\mathrm{exp}},
  \mathsf{Anc}_\zeta)
  :=
  \{\Phi_\eta:\eta\in\mathfrak S_{\pitchfork}\!\cdot d\}
  =
  \{\Phi_{-d},\Phi_0,\Phi_d\}.
\]
The exchange split is retained in $r_{\pitchfork}$ rather than scalarised a
second time as fragmentation, so this projection sets $\mu=0$. For arbitrary
$\lambda\geq0$, its Cogenesis persistence functional is
\begin{equation}
\label{eq:finite-riemann-persistence}
  \Persist_{X,T,\gamma,d}[\Phi_\eta]
  =
  -|r_X^{\mathrm H}(\eta)|^2
  +\lambda\rho_{\Phi_\eta}
  =
  \lambda-\eta^2.
\end{equation}
\end{definition}

\begin{proposition}[Unique finite persistence selection]
\label{prop:unique-riemann-persistence-selection}
At every finite declared resolution,
\[
  \Persist[\Phi_0]=\lambda,
  \qquad
  \Persist[\Phi_{\pm d}]=\lambda-d^2,
\]
and hence
\begin{equation}
\label{eq:unique-riemann-persistence-selection}
  \operatorname*{arg\,max}_{\Phi\in\Admissible}
  \Persist[\Phi]
  =\{\Phi_0\}.
\end{equation}
The ordering is independent of $X,T,\gamma,\lambda$, and attained magnitude.
\end{proposition}

\begin{proof}
Substitute \cref{eq:exact-branch-agreement} into
\cref{eq:finite-riemann-persistence}. Since $d>0$, the difference between the
central and either exchanged branch is exactly $d^2>0$.
\end{proof}

\begin{remark}[Scope of the finite selection]
\label{rem:finite-selection-scope}
\Cref{prop:unique-riemann-persistence-selection} selects the self-dual
\emph{resolution record}. It does not assert that
$\xi(\tfrac12+i\gamma)=0$, does not identify an unselected branch with its
projection, and does not by itself imply that an analytic zero must be the
selected unary. The latter follows only when the static zero equality is
carried through its complete causal ancestry, as isolated in
\cref{rev:foundational-inference}.
\end{remark}

\subsection{Complete per-zero normalisation record}

\begin{definition}[Per-zero normalisation record]
\label{def:per-zero-normalisation-record}
For a right-hand off-line candidate
$\rho=\tfrac12+d+i\gamma$ and finite $R\geq1$, define
\begin{equation}
\label{eq:per-zero-normalisation-record}
  \mathfrak B_{\rho;X,T,R}
  :=
  \left.
  \heartsim\!\left(\DeltaFork(\Field_\rho)\right)
  \right|_{(\Cctx_{X,T,\gamma,d},\Rres_{\mathrm{exp}})}
  =
  \odot_{\rho;X,T,R}^\ast
  +
  r_{\pitchfork,\rho;X,T,R}
  +\Delta'_{\rho;X,T,R}.
\end{equation}
Its components are
\[
 \begin{aligned}
  \odot_{\rho;X,T,R}^\ast&=\Phi_0,\\
  r_{\pitchfork,\rho;X,T,R}
  &=
  \left\{
    \Phi_{-d},\Phi_d,
    R_{\rho;X,T,R}^{-},R_{\rho;X,T,R}^{+},
    \widehat R_{\rho;X,T,R}
  \right\},\\
  \Delta'_{\rho;X,T,R}&=\tau|_{\{-d,d\}}.
 \end{aligned}
\]
Consequently its selected unary has
\[
  \pi_\delta(\odot_{\rho;X,T,R}^\ast)=0.
\]
The records $\Phi_{\pm d}$ retain their persistence values and ancestry; the
members $R_{\rho;X,T,R}^{\pm}$ and $\widehat R_{\rho;X,T,R}$ retain their
separately typed residual content and ancestry. The component
$\Delta'_{\rho;X,T,R}$ retains the exchange
$\Phi_d\leftrightarrow\Phi_{-d}$. The plus signs in
\cref{eq:per-zero-normalisation-record} denote typed co-retention, not scalar
addition.
\end{definition}

\begin{proposition}[Refinement-stable selection]
\label{prop:riemann-selection-refinement}
For all finite $X,T>1$, $R\geq1$, every $\gamma\in\mathbb R$, and every
$0<d<\tfrac12$, the selected coordinate of
\cref{eq:per-zero-normalisation-record} is
\begin{equation}
\label{eq:riemann-selection-refinement}
  \pi_\delta(\odot_{\rho;X,T,R}^\ast)=0.
\end{equation}
Changing any finite arithmetic, continuation, or transform cutoff transports
the residual ledger but does not change this selection.
\end{proposition}

\begin{proof}
The persistence difference is exactly $d^2$ by
\cref{prop:unique-riemann-persistence-selection} and contains no $X,T,\gamma$,
prime cutoff, continuation order, or endpoint. All such refinements are
transported by
\cref{prop:euler-cutoff-transport,prop:finite-continuation-cocycle,prop:theta-residual-transport},
while \cref{eq:exact-branch-agreement} remains unchanged. The selected
coordinate is therefore $0$ at every finite refinement.
\end{proof}

\begin{proposition}[No-orphan closure of the per-zero record]
\label{prop:riemann-no-orphan-closure}
Every identity, tail, cutoff schedule, phase, weighting, and residual appearing
in $\mathfrak B_{\rho;X,T,R}$ retains its dependency path to $\Field_\zeta$.
Neither sequential refinement nor the representation of $\Phi_{\pm d}$
introduces another persistence functional or normalising relation.
\end{proposition}

\begin{proof}
\Cref{prop:zeta-carrier-ancestry} establishes the dependency path for every
finite carrier, while \cref{lem:riemann-tail-stability} establishes the same
path for every finite tail of the exponent differential. In particular,
$\Pi_X^{\mathrm{HM}}$ is an explicit function of $s$ and $X$;
$\pi_{\mathrm{sd}}$, $\pi_{\mathrm{M,sd}}$,
$\widehat\pi_{\mathrm{sd}}$, $\pi_c$, and $\pi_\delta$ are declared
projections of that record; the scale residuals satisfy
\cref{eq:exact-branch-agreement}; and the finite and lifted projection
residuals $R_{\rho;X,T,R}^{\pm}$ and $\widehat R_{\rho;X,T,R}$ are retained by
\cref{def:finite-riemann-persistence,eq:lifted-self-dual-projection}. Moreover,
\cref{eq:riemann-branch-enactment} derives the finite family from the source
branch, its exchange, and its self-dual projection. The weights
$\lambda\geq0$ and $\mu=0$ are exactly those of the one declared persistence
functional in the declared scalar channel. The complete record
\cref{eq:per-zero-normalisation-record} applies the same causal normaliser
$\heartsim$ once and retains every unselected component. Projection changes
the expressed coordinate and residual, not the generator or normaliser.
\end{proof}

\begin{proposition}[Status of an off-line exponent candidate]
\label{prop:off-line-candidate-status}
For a hypothetical off-line zero, $\mathcal P_\rho$ is retained in
$r_{\pitchfork,\rho;X,T,R}$ or $\Delta'_{\rho;X,T,R}$ at every finite
refinement;
it is not the selected unary of the exponent-resolution channel.
\end{proposition}

\begin{proof}
\Cref{prop:riemann-differential-typing} types $\mathcal P_\rho$ as the
tail-stable candidate, and the finite projection $q_{X,T,R}^{+}$ in
\cref{def:finite-riemann-persistence} carries it into $\Phi_d$ with every
omitted component retained in $R_{\rho;X,T,R}^{+}$. The complete record retains
$\Phi_d$ explicitly, and
\cref{prop:unique-riemann-persistence-selection,prop:riemann-selection-refinement}
select $\Phi_0$ instead. Anti-collapse
\cref{eq:imported-projection,eq:imported-branch-normalisation} forbids replacing
either record by the other.
\end{proof}

\begin{reviewbox}
\textbf{What finite persistence proves.}
The scale residual is derived twice, the persistence comparison is exact, the
self-dual record is the unique finite maximum, and the paired off-line records
remain explicit residual content. The remaining foundational passage is not
that $\pm d$ disappear; it is that zero is the accounting projection of a
causally completed zero-forming process, not a primitive point that exists
before that process. The explicit chain
\cref{eq:imported-finite-selection,eq:imported-branch-normalisation} types its
realised output as the selected unary and its unselected records as residual
or continuing content.
\end{reviewbox}

\section{Processual proof contract}
\label{sec:proof-contract}

\begin{keybox}[title={\bfseries Processual proof contract},fontupper=\small]
\begin{enumerate}[label=\textbf{\arabic*.},leftmargin=*]
  \item Declare every domain, context, resolution, projection, and residual.
  \item Prove the reformulation exactly equivalent to RH.
  \item Derive the finite process from native ancestry.
  \item Recover static symmetry from the process, not conversely.
  \item Prove magnitude independence rather than naming it.
  \item Prove invariant preservation at every finite refinement.
  \item Translate an off-line zero into \cref{eq:off-line-zero-differential}.
  \item Exhaust every carrier without presuming exclusion.
  \item Prove \cref{eq:required-local-implication}.
  \item Isolate the foundational causal extension at
        \cref{rev:foundational-inference}.
  \item Recover \cref{eq:riemann-target} by ordinary complex arithmetic.
  \item Map every falsifiable claim to a verifier that fails non-zero.
\end{enumerate}
\end{keybox}

\section{The single foundational causal extension}
\label{sec:foundational-inference}

\begin{remark}[The single foundational causal extension]
\label{rev:foundational-inference}
Let
$\rho=\tfrac12+\delta+i\gamma$ be the right-hand representative of a
hypothetical off-line orbit. Writing $\xi(\rho)=0$ as the claimed scalar
output of its complete record, the decisive causal order is:
\begin{equation}
\label{eq:foundational-inference}
\begin{aligned}
  \mathsf{Anc}_\rho
  &\longrightarrow
  (\Cctx_\rho,\Rres_\rho)
  \longrightarrow
  \Pi_\rho
  \longrightarrow
  \Persist_\rho
  \longrightarrow
  \odot_\rho^\ast+R_\rho,\\
  \odot_\rho^\ast+R_\rho
  &\overset{\pi_\xi}{\longrightarrow}
  \underbrace{
    \mathcal Z_\rho
    =\xi(\rho)+r_{0,\rho}
    =0+r_{0,\rho}
  }_{\text{static zero accounting projection}},\\
  &\Longrightarrow
  \underbrace{
    \Phi_\delta=\odot_\rho^\ast
  }_{\text{claimed output occupies the selected unary}}
  \Longrightarrow
  \underbrace{
    \pi_\delta(\Phi_\delta)
    =\pi_\delta(\odot_\rho^\ast)=0
  }_{\delta=0}.
\end{aligned}
\end{equation}
The domain relations in this chain are finite and exact:
\cref{prop:exact-branch-agreement} identifies the spectral residual with the
factor-exchange residual,
\cref{prop:unique-riemann-persistence-selection} gives the unique selected
record, and \cref{prop:riemann-selection-refinement} preserves that selection
at every finite refinement. The paired records $\Phi_{\pm\delta}$ are not
deleted; \cref{prop:off-line-candidate-status} retains them as residual or
continuing content.

The implication following the terminal zero projection is causally explicit,
not an additional selection axiom. Claiming the equality $\xi(\rho)=0$ as a
realised zero-forming output means that $\rho$ names the carried output of the
ancestry-first process displayed above; it does not make that static equality
the process input.
The complete normalisation relation
\cref{eq:imported-finite-selection,eq:imported-branch-normalisation} types the
carried output as $\odot_\rho^\ast$ and the other expressions as residual or
continuing content. To call $\Phi_\delta$ both an unselected residual and the
realised output in the same channel would reverse that explicit causal order.
Carrier closure \cref{prop:riemann-no-orphan-closure} ensures that refinement
introduces no second normaliser under which its status could change.

A completed-infinite reading can represent an off-line coordinate as a member
of a static completed zero set. That representation supplies no competing
zero-forming ancestry: it promotes the accounting projection to a primitive
object and terminates causal accounting at the equality. Sequential causal
completion instead retains that coordinate as residual while making only the
normalised unary available as the zero-forming output. This is the same
selected-versus-continuing causal distinction used in the Collatz orbit
argument.
\end{remark}

\begin{theorem}[Critical-line closure]
\label{thm:critical-line-closure}
For every $\rho\in\CriticalStrip$ and arbitrary finite $X,T>1$, $R\geq1$,
\begin{equation}
\label{eq:required-local-implication}
  \xi(\rho)=0\ \land\ \rho\neq\Dual(\rho)
  \Longrightarrow
  \mathcal P_{\rho^+}
  \land\Audit_\zeta(\rho^+)
  \land\Phi_\delta\in r_{\pitchfork,\rho^+;X,T,R}
  \Longrightarrow
  \delta=0
  \quad\bot.
\end{equation}
Here $\rho^+\in\{\rho,\Dual(\rho)\}$ is the right-hand representative
$\rho^+=\tfrac12+\delta+i\gamma$ with $\delta>0$.
\end{theorem}

\begin{proof}
If the zero is off-line, \cref{cor:off-line-zero-implies-persistence} gives
$\mathcal P_{\rho^+}$. The derivation and
\cref{prop:zeta-carrier-ancestry,prop:riemann-no-orphan-closure} establish
$\Audit_\zeta(\rho^+)$ and the complete record
$\mathfrak B_{\rho^+;X,T,R}$. The candidate has spectral coordinate
$\pi_\delta(\Phi_\delta)=\delta$, while finite persistence gives
$\pi_\delta(\odot_{\rho^+;X,T,R}^\ast)=0$ at every refinement. The foundational
causal extension \cref{eq:foundational-inference} carries the static zero
projection through its complete ancestry to the latter selected coordinate,
hence $\delta=0$, contradicting the right-hand orientation. Therefore no
off-line orbit is realised as a zero-forming output, and every realised zero
projection is fixed by $\Dual$. The records $\Phi_{\pm\delta}$ remain retained
as residual candidates by \cref{prop:off-line-candidate-status}; they are not
erased. Since $\Dual$ is involutive, the zero-forming conclusion applies to
the originally named member $\rho$ as well as $\rho^+$.
\end{proof}

\section{Main theorem and classical recovery}
\label{sec:classical-recovery}

\begin{theorem}[Riemann hypothesis]
\label{thm:riemann-hypothesis}
Every nontrivial zero $\rho$ of $\zeta$ satisfies
\[
  \Re(\rho)=\frac12.
\]
\end{theorem}

\begin{proof}
Inside $\CriticalStrip$, the factors multiplying $\zeta(s)$ in
\cref{eq:completed-xi} are nonzero, so the nontrivial zeros of $\zeta$ and
$\xi$ coincide. Critical-line closure
\cref{thm:critical-line-closure} gives $\rho=\Dual(\rho)$. By
\cref{prop:fixed-locus}, $\Fix(\Dual)=\CriticalLine$, and therefore
$\Re(\rho)=\tfrac12$.
\end{proof}

\section{Conventional-language proof skeleton}
\label{sec:conventional-skeleton}

\begin{enumerate}[label=\textbf{\arabic*.},leftmargin=*]
  \item Unique factorisation generates the finite convolution algebra for
        Dirichlet coefficients and its compatible Möbius inverse.
  \item Exact scalarisation, cutoff, Bernoulli, theta, and Mellin residuals
        transport that ancestry without identifying a projection with its
        generator.
  \item The finite Möbius--hyperbola ledger is identically $1$ at every
        resolution; factor exchange has the unique fixed exponent
        $\theta=\tfrac12$.
  \item Theta--Mellin transport projects reciprocal scale to
        $s\leftrightarrow1-s$, recovering the functional equation rather than
        assuming it.
  \item An off-line zero
        $\rho=\tfrac12+\delta+i\gamma$ forces the exact persistent Möbius
        obstruction $\mathcal P_\rho$, an $X^\delta$ differential relative to
        square-root resolution.
  \item Under the declared spectral--hyperbola projection, the hyperbola and
        Mellin records carry the same exponent residual $\delta$. The Cogenesis
        finite persistence functional has the
        unique maximum $\delta=0$ and retains $\pm\delta$ as unselected
        branches.
  \item The per-zero normalisation record preserves that selection through
        every finite refinement while carrier closure retains every ancestry
        and residual.
  \item Extending causal accounting through the complete zero-forming record
        makes the selected unary its realised output and retains every
        unselected branch as residual, so $\delta=0$.
  \item The fixed locus of $1-\overline{s}$ is the critical line.
\end{enumerate}
The eighth step is the foundational causal extension isolated in
\cref{rev:foundational-inference}.

\section{Verification architecture}
\label{sec:verification}

Verification follows the Collatz precedent \parencite{petrik2026collatz}:
scripts exit non-zero on failure and are audited against theorem labels.
The executable clusters are:
\begin{itemize}
  \item \path{verification/verify_finite_arithmetic.py}: prime ancestry, restriction,
        Möbius and Euler inversion, scalar boundaries, cutoff residuals, and
        additive--multiplicative ledgers;
  \item \path{verification/verify_continuation_transport.py}: exact rational
        Euler--Maclaurin cocycles, finite theta transport, the
        Gaussian--Poisson bridge, and finite Mellin symmetry;
  \item \path{verification/verify_hyperbola_invariant.py}: arbitrary hyperbola cutoffs,
        inverse-convolution widening, self-dual resolution, and the finite
        Möbius--Mellin ledger;
  \item \path{verification/verify_persistence_selection.py}: exact
        cross-projection exponent agreement, widened registration resolution,
        finite persistence ordering, and $X$-, $T$-, and height-stable
        selection; and
  \item \path{verification/verify_claim_coverage.py}: every labelled theorem,
        proposition, and corollary is mapped to an executable verifier or an
        explicit analytic or causal-process review obligation.
\end{itemize}
The aggregate runner \path{verification/run_all.py} fails on the first
unsatisfied claim. The current suite performs $18377$ checks across $35$
labelled claims: $22$ executable and $13$ analytic or causal-process-level. Code
verifies the finite domain claims around the isolated causal extension and
audits that no projection, residual, or ancestry has been hidden inside it.
The causal order itself is explicit in the typed process record rather than an
empirical claim delegated to computation.

\section{Reader guardrails}
\label{sec:guardrails}

\begin{longtable}{@{}p{0.42\textwidth}p{0.52\textwidth}@{}}
\caption{Common misreadings and required corrections.}
\label{tab:guardrails}\\
\toprule
\textbf{Misreading} & \textbf{Required correction} \\
\midrule
\endfirsthead
\toprule
\textbf{Misreading} & \textbf{Required correction} \\
\midrule
\endhead
Functional symmetry fixes every zero. &
It closes the zero set under reflection; off-line orbits remain compatible. \\
\addlinespace
A zero is a primitive point that exists before its generating relation. &
Zero is the accounting projection of completed cancellation in a declared
channel; its ancestry, residual, and normalisation remain causally prior. \\
\addlinespace
The fixed locus is the critical line, therefore RH follows. &
The missing implication is that every zero must be fixed. \\
\addlinespace
Infinity is not completed, therefore no off-line zero exists. &
An off-line zero is finite and requires an exact local exclusion. \\
\addlinespace
A numerical pattern is the invariant. &
Numerics may falsify but cannot derive or universally establish it. \\
\addlinespace
A classical representation is automatically the generator. &
It is a projection whose causal role follows from its exact dependency. \\
\addlinespace
A residual tends to zero, so it may be discarded. &
$r\to0$ is not $r=0$; exact zero must be proved in the declared register. \\
\addlinespace
Carrier exhaustion excludes internally generated counterexamples. &
It classifies their support. The finite persistence comparison and complete
normalisation record are still required. \\
\addlinespace
The finite persistence maximum alone proves RH. &
It selects a self-dual resolution record. Zero becomes a realised output only
after the complete causal chain reaches that selected unary; the other records
remain explicitly residual. \\
\bottomrule
\end{longtable}

\section{Discussion}
\label{sec:discussion}

The Collatz analogy is methodological, not evidentiary. No Collatz quantity,
map, valuation, or conclusion is imported. The proof desymmetrises the
functional equation by recovering the finite process it projects, then derives
a preserved relation not assumed from RH. The relation is widened beyond the
target case so computation can falsify it rather than merely confirm fitted
behaviour. The persistence layer instantiates the same selected-versus-
continuing distinction used in the Collatz orbit proof: the domain quantities
change, but the explicit extension of causal availability through
normalisation does not.

\section{Conclusion}
\label{sec:conclusion}

The proof begins with relation rather than static symmetry, derives the
height-independent arithmetic invariant and its square-root resolution,
retains finite residual transport, types an off-line zero as an exact
$X^\delta$ Möbius obstruction, derives the same exponent residual in hyperbola
and Mellin projections, and computes the self-dual record as the unique finite
persistence selection. Paired off-line records remain explicit residual
content. The isolated foundational causal extension carries the zero accounting
projection through its complete ancestry to the selected unary; its fixed
locus is the critical line, yielding \cref{thm:riemann-hypothesis}.

\appendix
\section{Derivation ledger}
\label{app:derivation-ledger}

Every substantive revision must record the claim, native ancestry, discovery
or validation status, exact derivation, projection boundary, residual,
verifier, reason for change, and revert path.

\section{Submission and verification notes}
\label{app:submission}

\begin{itemize}
  \item Compile with
        \texttt{latexmk -pdf -interaction=nonstopmode -halt-on-error}.
  \item Maintain citations in \texttt{references.bib}.
  \item Supply verification scripts beside the manuscript when claims exist.
  \item Record causal dependencies at
        \cref{sec:cogenesis-dependencies} and the proof derivation at
        \cref{sec:native-ancestry,sec:process-invariant,sec:counterexample-differential,sec:foundational-inference}.
\end{itemize}

\printbibliography[title={References}]
\end{document}
